\documentclass[11pt]{article}
\usepackage[T1]{fontenc}\usepackage{lmodern}\usepackage[sc]{mathpazo}
\usepackage{amsmath,amssymb,amsthm,mathtools,booktabs,array,geometry,fancyhdr,xcolor}
\geometry{paperwidth=7.65in,paperheight=10.4in,inner=.83in,outer=.78in,top=.80in,bottom=.85in}
\usepackage[unicode,pdfencoding=auto]{hyperref}
\definecolor{teal}{HTML}{20555A}\hypersetup{colorlinks=true,linkcolor=teal,urlcolor=teal,citecolor=teal,pdftitle={Averaged zero-free strips for automorphic L-functions}}
\pagestyle{fancy}\fancyhf{}\fancyhead[L]{\small\itshape Averaged zero-free strips}\fancyhead[R]{\small\thepage}
\setlength{\headheight}{15pt}\setlength{\emergencystretch}{2em}\setcounter{tocdepth}{1}
\newtheorem{theorem}{Theorem}[section]\newtheorem{lemma}[theorem]{Lemma}\newtheorem{corollary}[theorem]{Corollary}
\newtheorem{proposition}[theorem]{Proposition}\theoremstyle{definition}\newtheorem{definition}[theorem]{Definition}\newtheorem{remark}[theorem]{Remark}
\newcommand{\R}{\mathbb R}\newcommand{\Q}{\mathbb Q}\newcommand{\C}{\mathbb C}\newcommand{\Z}{\mathbb Z}
\newcommand{\GL}{\mathrm{GL}}\newcommand{\Sym}{\mathrm{Sym}}\newcommand{\re}{\operatorname{Re}}\newcommand{\im}{\operatorname{Im}}
\newcommand{\abs}[1]{\lvert #1\rvert}\newcommand{\ord}{\operatorname{ord}}
\begin{document}
\begin{center}
{\color{teal}\sffamily MATHEMATICS PROJECT}\par\medskip
{\LARGE Averaged zero-free strips for automorphic $L$-functions}\par\medskip
{\large An explicit amplifier for elliptic curves}\par\smallskip
{\small Research supplement \quad $\cdot$ \quad 23 November 2026}
\end{center}
\begin{abstract}
For each fixed non-CM elliptic curve over $\Q$ we derive an asymptotic zero-free strip with denominator $5.7693\log|T|$ in unitary normalization. The construction uses a common translation in the ordinates of finitely many symmetric-power logarithmic derivatives. A polynomial on the square of Frobenius and Fourier parameters controls all good prime powers simultaneously. One certificate uses 170 terminal rectangles to control the symmetric square. A second uses 376 to bound zeros of the original curve. We give the proof of the transfer from the weighted prime sum to the zero of the original curve and identify the external arithmetic results it uses. Separate arguments treat general automorphic representations, simultaneous zeros, and an effective strip for the curve of conductor eleven. The final sections identify the part of each argument implemented in Lean 4. The extended Riemann hypothesis over number fields remains open.
\end{abstract}
\makeatletter
\renewcommand*\l@section{\@dottedtocline{1}{0em}{2.1em}}
\makeatother
{\small\tableofcontents}
\clearpage
\section{Statement and normalization}
Write $L_1(s)=L(E,s+1/2)$, where the right-hand side is the Hasse--Weil function in arithmetic normalization. Write $L_j(s)$ for the unitarily normalized $j$-th symmetric power and $L_0(s)=\zeta(s)$. The nontrivial zeros here are those of the completed function in $0<\re s<1$.
\begin{theorem}\label{main}
For every fixed non-CM elliptic curve $E/\Q$ there exists $T_E>1$ such that for every nontrivial zero $\beta+iT$ of $L_1$ with $|T|\ge T_E$,
\[
       \min(\beta,1-\beta)>\frac{1}{5.7693\log|T|}.
\]
Equivalently, in arithmetic notation all such zeros lie strictly between $1/2+1/(5.7693\log |T|)$ and $3/2-1/(5.7693\log |T|)$ in real part.
\end{theorem}
The constant in the preceding manuscript was $12.326$ for this case. The new denominator is a fixed-height-asymptotic result for each curve. Section~\ref{section:level11-effective} gives a different constant with a specified starting height. Both assertions leave open the location of zeros in the intervening strip. The first quantitative step is the positive-average denominator $6.16918$. Section~\ref{sec:final} uses the symmetric square to improve it to $5.7693$. The theorem is a deduction from published automorphy, subconvexity and classical zeta bounds. Its new ingredients within this project are the choice of signed symmetric-power coefficients, their positivity certificate and their combination with a compact averaging test. We make no priority comparison with every choice of test in the literature.
\section{A disk estimate with one retained zero}
For a finite $L$-function $F$, put $D_F(s)=-\re(F'/F)(s)$. Suppose $F$ has no zero to the right of $1$, has the usual finite-order strip growth, and
\[
 |F(1/2+it)|\ll_{F,\epsilon}(1+|t|)^{\gamma+\epsilon},
 \qquad |\log |F(1+d+it)||\ll_F\log(2+1/d)\quad(0<d\le1).
\]
All disks below are drawn at large ordinate and are free of poles. This version of the disk method follows Heath-Brown, Lemma 3.2 \cite{HB}.
\begin{lemma}\label{disk}
Fix $a,v>0$, a positive integer $\ell$, and $M<\infty$. Put $L=\log|T|$, $\delta=a/L$, $s=1+v\delta+i(\ell T+\delta\lambda)$, $|\lambda|\le M$. There is a radius $1/4\le R\le1/3$ avoiding boundary zeros for which
\[
 D_F(s)\le (\gamma+o(1))L-
 \sum_{|s-\rho|<R}(\re s-\re\rho)
                  \big(|s-\rho|^{-2}-R^{-2}\big).
\]
Here and below an upper bound with $\gamma+o(1)$ means a limsup at most $\gamma$ after normalization by $L$. The error is uniform over these $\lambda$ and radii. Each summand in the displayed sum is nonnegative.
\end{lemma}
\begin{proof}
Choose any allowed radius avoiding zeros on its circumference. In local coordinate $z$ about $s$ remove the zeros $w=\rho-s$, with multiplicity, using the disk Blaschke factors $R(z-w)/(R^2-\bar w z)$. Their moduli on $|z|=R$ equal one. The logarithm of the modulus of the quotient is harmonic. Its horizontal derivative at the center, by the Poisson integral, is $(\pi R)^{-1}\int_0^{2\pi}\log |F(s+Re^{i\theta})|\cos\theta\,d\theta$. The real logarithmic derivative of an individual factor at zero is $(\re s-\re\rho)(|s-\rho|^{-2}-R^{-2})$. This proves the identity with the boundary integral taken with a minus sign in $D_F$.
 
On $\cos\theta\ge0$ the Euler estimate gives a lower bound $-O(\log L)$. On $\cos\theta<0$, apply Phragm\'en--Lindel\"of between $1/2$ and $1+\epsilon$ with fixed $\epsilon>0$. The upper bound, divided by $L$, is at most $2\gamma(-R\cos\theta)+O_F(\epsilon)+o_{F,\epsilon}(1)$. To see the slope, linearly interpolate the two boundary exponents $\gamma+\epsilon$ and zero over a width $1/2+\epsilon$ and use $\log(2+|\ell T+O(1)|)=L+O(1)$. Beyond the right boundary the Euler product is bounded at fixed $\epsilon$. The integral of $\cos^2\theta$ over this left semicircle is $\pi/2$. First let $|T|$ grow, then let $\epsilon$ decrease to zero. Since $R\ge1/4$, the bound is uniform.
\end{proof}
\section{Averaging at the ordinate of the original zero}
Let $T_\ell(\cos\theta)=\cos(\ell\theta)$ and $U_j(\cos\theta)=\sum_{k=0}^j e^{i(j-2k)\theta}$. At the endpoints use the polynomial continuations. For a finite table $(w_{\ell j})$ write
\[
  A(x,y)=\sum_{\ell,j\ge0} w_{\ell j}U_j(x)T_\ell(y).
\]
For the moment use $L_1$ as the target. The proofs also apply with any designated mixed index $(1,j_0)$ of weight one, retaining the corresponding zero of $L_{j_0}$. Suppose $A\ge0$ on $[-1,1]^2$, $w_{11}=1$, $P=w_{00}\ge0$, and $w_{\ell j}\ge0$ whenever $\ell j>0$. The other coefficients can have either sign. Put $G=\sum_{\ell j>0}w_{\ell j}\gamma_j$, where $\gamma_j$ is a central-line exponent in Lemma~\ref{disk}. Let $\mu$ be an even probability measure on $\R$ whose cosine transform $h_b(y)=\int\cos(\lambda y)d\mu(\lambda)$ is nonnegative on the real line and integrable. Put
\[
 H(z)=\int_0^\infty h_b(y)e^{-zy}\,dy\quad(z\ge0).
\]
\begin{proposition}\label{average}
Suppose the finite functions used in the table satisfy the hypotheses of the disk lemma at moving positive-degree indices, the static $L_j$ with $j>0$ are holomorphic and nonzero at $1$, and
\[
\begin{split}
 &D_{L_0}(1+d+i\tau)=o(\log |\tau|)\quad
       \text{as }|\tau|\to\infty,\quad d\asymp1/\log|\tau|,\\
 &|D_{L_j}(1+d+i\tau)|\ll_j 1/d \qquad(0<d\le1,\ \tau\in\R).
\end{split}
\]
At good primes suppose the Euler traces of $L_j$ at $p^\nu$ are $U_j(\cos(\nu\theta_p))$ with real $\theta_p$. Suppose the total from the remaining finite set of primes is bounded at $\re s\ge1$. Then every fixed $a>0$ satisfying $aG<H(1)-PH(0)$ has the following property. At all sufficiently large ordinates every zero $\beta+iT$ of $L_1$ satisfies $(1-\beta)\log|T|>a$.
\end{proposition}
\begin{proof}
Consider a contrary sequence $|T|\to\infty$. Keep the notation $\delta=a/L$ and put $w=(1-\beta)/\delta$. By passing to a subsequence we may assume $w\to w_*\in[0,1]$. At each index evaluate $D_{L_j}$ at $s_\ell=1+v\delta+i(\ell T+\delta\lambda)$ for fixed $v>0$. The displacement $\delta\lambda$ is common to all indices, including $\ell=0$. In the integrated Euler sum the summand at good $p^\nu$, apart from the factor $(\log p)p^{-\nu(1+v\delta)}$, is
\[
 h_b(\delta\nu\log p)\ A\big(\cos(\nu\theta_p),\cos(T\nu\log p)\big).
\]
Indeed the integrated sine contribution vanishes by evenness. Absolute convergence at each evaluation abscissa permits interchanging the logarithmic derivative series and the probability average. Thus the full weighted integral is bounded below by $-O_E(1)$.

On $|\lambda|\le M$ apply the disk lemma to all moving terms with $j>0$. Their coefficients are nonnegative. In the $(1,1)$ term retain the selected zero once, and discard all other zero contributions. For large $T$ this zero belongs to every admissible target disk, since its distance from the new center is $\delta\sqrt{(v+w)^2+\lambda^2}$. After multiplication by $\delta$, its contribution equals
\[
 \frac{v+w}{(v+w)^2+\lambda^2}
            -\frac{\delta^2(v+w)}{R^2}.
\]
The second term tends uniformly to zero on this compact range. The static $j>0$ terms contribute $o(1)$ after scaling, even with negative coefficients. The static zeta pole gives $P v/(v^2+\lambda^2)+o(1)$. All moving zeta terms contribute $o(1)$ by the scalar assumption, for either sign. On the complement of this compact interval use the absolute Euler bound on the original, unexpanded integrand. Its contribution after scaling is at most $C\mu(|\lambda|>M)/v$ in absolute value, independently of $T$. This tail argument permits the disk radius and the cutoff to be chosen without any estimates for zeros at remote translates.

Let $|T|\to\infty$ first and $M\to\infty$ second. For $z>0$, the identity
\[
 \int\frac z{z^2+\lambda^2}\,d\mu(\lambda)
    =\int_0^\infty e^{-zy}\int\cos(\lambda y)d\mu(\lambda)\,dy=H(z)
\]
follows by integrating $e^{-zy}\cos(\lambda y)$ and applying Fubini with majorant $e^{-zy}$. Hence $aG\ge H(v+w_*)-PH(v)\ge H(v+1)-PH(v)$. The last inequality uses $h_b\ge0$. Dominated convergence as $v\downarrow0$ gives the contradiction.
\end{proof}
\paragraph{Application of the analytic inputs.}
By modularity and Newton--Thorne, $L_j$, $j>0$, is the standard function of a fixed unitary cuspidal representation on $\GL_{j+1}/\Q$ \cite{BCDT,NT}. Its standard continuation, finite-order strip bounds and nonvanishing on $\re s=1$ apply \cite{JS}. Hasse's bound at good primes gives the angles. At each bad prime the normalized local Weil--Deligne parameter of an elliptic curve is tempered. For potentially good reduction it becomes unramified and pure after a finite extension. For potentially multiplicative reduction it is a finite character twist of a special parameter. In both cases the reciprocal roots in the inertia-invariant kernel of monodromy for a fixed symmetric power have absolute value at most one in unitary normalization. Its logarithmic derivative at $\sigma\ge1$ therefore costs at most $(j+1)\log p/(p-1)$. Deleting or reinstating the finitely many bad primes costs $O_{E,j}(1)$ uniformly in height. At the other primes $|U_j|\le j+1$, so the absolute Euler logarithm is bounded by $(j+1)\log\zeta(1+d)$ up to a constant and the absolute logarithmic derivative by $(j+1)(-\zeta'/\zeta)(1+d)+O(1)\ll_j1/d$. For the moving $j=0$ terms use the classical Vinogradov--Korobov estimate $\zeta'/\zeta(\sigma+i\tau)\ll(\log|\tau|)^{2/3}(\log\log|\tau|)^{O(1)}$ at $\sigma\ge1$ near one \cite{Titch}. These assertions apply to both signs of $T$ and are uniform under the bounded shifts used in the proof.
 
We use the exponent $\gamma_1=1/3$ of Booker--Milinovich--Ng for primitive holomorphic forms of arbitrary fixed level and, for $j\ge2$, the limiting exponent in Nelson's Theorem 1.1 \cite{BMN,Nelson},
\[
 \gamma_j=\frac{k}{4}\left(1-\frac{2}{3k^5-2k^4-k^2}\right),\qquad k=j+1 .
\]
Each exponent permits an arbitrarily small positive loss as required in Lemma~\ref{disk}. Static indices carry no such cost.


\section{A compact test with exact moments}
Put $f(u)=u(1-u)(1+\tfrac52u(1-u))$ on $[0,1]$ and extend it by zero to a continuous function $g$ on the real line. Direct integration gives $m=\int g^2=199/2520$. For $b>0$ take the positive scale change of the normalized autocorrelation
\[
  h_b(t)=\frac1m\int_{\R}g(u)g(u+t/b)\,du .
\]
The function is even and nonnegative and vanishes for $|t|\ge b$. With angular-frequency convention $\widehat g(\xi)=\int g(u)e^{-i\xi u}du$, the measure
\[
 d\mu_b(\lambda)=\frac{b}{2\pi m}|\widehat g(b\lambda)|^2d\lambda
\]
has total mass one and cosine transform $h_b$. Indeed Plancherel gives the mass, convolution gives the squared modulus, and Fourier inversion applies because the autocorrelation is continuous and its transform is integrable.

For $0\le x\le1$ the overlapping supports give
\[
\begin{split}
 h_b(bx)=h(x)&=\frac{2520}{199}\int_0^{1-x} f(u) f(u+x)\,du\\
 &=\frac{(1-x)^3}{199}\big(25x^6+75x^5-120x^4-560x^3\\[-3pt]
 &\hspace{40mm}+204x^2+597x+199\big).
\end{split}
\]
This identity proves the sign of $h$ directly from the integral. Substituting $t=bx$ gives $H(1)-PH(0)=b\int_0^1 h(x)(e^{-bx}-P)\,dx$.
\begin{lemma}
For $b=13/8$ the last quantity is at least the rational affine function
\[
 J_0(P)=\frac{13}{8}\left(\frac{95181399549361514143}{344366740312739020800}
                          -\frac{315}{796}P\right).
\]
\end{lemma}
\begin{proof}
Put $t_9(z)=\sum_{j=0}^9(-z)^j/j!$. Differentiation gives $(e^z t_9(z))'=-e^z z^9/9!$, so $t_9(z)\le e^{-z}$ on $z\ge0$. Integrate this inequality against $h\ge0$. Write $h(x)=\sum_{i=0}^9 h_i x^i$ by multiplying the displayed polynomial. Then
\[
 \int_0^1 h(x)t_9(bx)\,dx=\sum_{i=0}^9\sum_{j=0}^9\frac{h_i(-b)^j}{j!(i+j+1)},
 \qquad \int_0^1 h(x)\,dx=\sum_{i=0}^9\frac{h_i}{i+1}=\frac{315}{796}.
\]
Substituting $13/8$ in the first sum yields the stated fraction. The formal proof differentiates explicit rational-polynomial primitives and applies the fundamental theorem of calculus to the Lebesgue interval integral.
\end{proof}


\section{An octic amplifier for the smoothed inequality}
Write $c_{\ell j}$ for the entries in the following table and set
$A(x,y)=10^{-6}\sum c_{\ell j}T_\ell(y)U_j(x)$.
An omitted row or column has weight zero.
\[
\begin{array}{c|rrrrrrrr}
\ell\backslash j&0&1&2&3&4&5&6&8\\\hline
0&470377&0&412947&0&30575&0&2364&-675\\
1&0&1000000&0&247736&0&0&0&0\\
2&460679&0&338826&0&0&0&1821&0\\
3&0&193459&0&2514&0&0&0&0\\
4&37617&0&0&0&0&0&0&203\\
5&0&0&0&0&0&204&0&0\\
6&2046&0&888&0&0&0&0&0\\
8&-339&0&0&0&115&0&0&0
\end{array}
\]
\paragraph{Proposition.}
$A$ is nonnegative on $[-1,1]^2$. Its coefficient of $T_1U_1$ equals one. The only negative entries have either $\ell=0$ or $j=0$.
Give terms with $\ell j=0$ cost zero and set
\[
d_1=\frac13,\qquad d_j=\frac{j+1}{4}\left(1-\frac2{3(j+1)^5-2(j+1)^4-(j+1)^2}\right)\quad(j\ge2).
\]
For $P=c_{00}/10^6$ and $G=10^{-6}\sum_{\ell,j\ge1}d_j c_{\ell j}$, the quartic autocorrelation gain
\[
J(P)=\frac{13}{8}\left(\frac{95181399549361514143}{344366740312739020800}-\frac{315}{796}P\right)
\]
is positive and $G/J(P)<6.17635$.
\paragraph{Proof.}
Homogenize the integer numerator by putting $x=(a-b)/(a+b)$ and $y=(u-v)/(u+v)$, and multiply by $(a+b)^8(u+v)^8$.
For a homogeneous row $\sum_i r_i u^iv^{8-i}$ the coefficient of $s^kt^{8-k}$ after substituting $(u,v)=(s,s+2t)$ is
$2^{8-k}\sum_{i=0}^k\binom{8-i}{8-k}r_i$.
For the substitution $(u,v)=(2s+t,t)$ it is
$2^k\sum_{i=k}^8\binom{i}{k}r_i$.
The first chart covers $u\le v$ and the second covers $u\ge v$. Start with the first variable pair $(a,b)$, alternate the pairs, and stop as soon as all entries of a rectangle are nonnegative. The terminal counts at respective depths $3$ through $13$ are
\[
2,\ 2,\ 2,\ 8,\ 14,\ 34,\ 22,\ 62,\ 86,\ 156,\ 96 .
\]
There are $484$ terminal rectangles. A finite induction up the tree proves nonnegativity on the original square. The formal proof verifies every integral substitution and the exact identity between the root polynomial and the displayed coefficient table.
Direct summation gives
\[
P=\frac{470377}{10^6},\qquad G=\frac{85398805176211663}{94275997200000000},\qquad
J(P)=\frac{50505659116627493840711}{344366740312739020800000}.
\]
Consequently
\[
\frac{G}{J(P)}
=\frac{1473490268014203275091541753856}{238569794044779545593328497375}<\frac{617635}{100000}.
\]
The corresponding formal results are in \texttt{EllipticSmoothOctic.lean}.
\paragraph{Arithmetic scope.}
At good primes of a rational elliptic curve the normalized traces of all powers are in $[-1,1]$ by Hasse's bound. At a bad prime its contribution must instead be estimated in absolute value, uniformly in the ordinate. In using the explicit formula, the signed terms with $\ell=0$ and $j=0$ require the corresponding static and scalar bounds before discarding the other zero contributions. These estimates belong to the analytic reduction.


\section{A degree-twenty amplifier}
The same quartic autocorrelation test admits a smaller reciprocal width coefficient. Use the normalization $A=10^{-6}\sum_{\ell,j}c_{\ell j}T_\ell U_j$ and the costs $G$ and $J$ of the octic construction. In the table an entry $(j,c)$ in row $\ell$ means $c_{\ell j}=c$. Unlisted entries vanish.
\[
\begin{array}{c|l}
\ell &(j,c)\\ \hline
0 &(0,473580),(2,408146),(4,28812),(6,2480),(8,-204),(10,-1223)\\
  &(12,761),(14,-23),(16,-181),(18,84),(20,2)\\
1 &(1,1000000),(3,239816)\\
2 &(0,465965),(2,335894)\\
3 &(1,195965)\\
4 &(0,39244)\\
5 &(5,828)\\
6 &(0,1962),(2,1676)\\
8 &(0,77)\\
9 &(1,585)\\
10 &(0,44)\\
12 &(0,198)\\
13 &(1,102)\\
14 &(0,-62)\\
16 &(0,104),(2,43)\\
18 &(0,-18)\\
19 &(1,71)\\
20 &(0,48)
\end{array}
\]
\paragraph{Proposition.}
$A(x,y)\ge0$ for $|x|, |y|\le1$, all its moving nonscalar coefficients are nonnegative, and
\[
   \frac{G}{J(P)}=\frac{63091694191815282764267126784}{10226927578976034742586005625}
   <\frac{616918}{100000}.
\]
\paragraph{Proof.}
The nonnegativity follows from the same two integer substitutions as in the octic construction with the degree changed from $8$ to $20$. For reproducibility, the bidegree-$(20,20)$ root coefficient matrix is
\[
 H_{ik}=\sum_{\ell,j}c_{\ell j}\,[z^i](1+z)^{20}
 U_j\bigl((z-1)/(z+1)\bigr)\,[w^k](1+w)^{20}
 T_\ell\bigl((w-1)/(w+1)\bigr).
\]
Starting with the first pair, alternating axes and bisecting only rectangles whose matrix has a negative entry, the terminal counts at depths $3$ through $12$ are
\[
 2,\ 2,\ 4,\ 8,\ 14,\ 28,\ 34,\ 64,\ 46,\ 20.
\]
Each of the $222$ terminal matrices is entrywise nonnegative. The two children cover their parent at every bisection, which proves the claim at the root. The Lean theorem \texttt{SmoothDegree20.normalized\_fourier\_nonnegative} proves the resulting real inequality. The standard-library verifier \texttt{check\_smoothing.py} computes the displayed root matrix from the Chebyshev recurrences, performs all substitutions with arbitrary-precision integers and checks the slope and gain by exact rational arithmetic. They equal
\[
 P=473580/10^6,\qquad
 G=\frac{219852918901}{246450000000},\qquad
 J(P)=\frac{9959272140207946188763}{68873348062547804160000}.
\]
Division gives the asserted upper bound. Independently, \texttt{EllipticSmoothDegree20.lean} proves the same normalized Fourier nonnegativity in Lean. It checks bisection on $x\le0$ and obtains the other half from $A(-x,-y)=A(x,y)$, since each occupied pair has $j\equiv\ell\pmod2$. Its theorem \texttt{SmoothDegree20.ratio\_lt} proves the strict rational bound. The analytic estimates needed for the signed static and scalar terms have the same scope as in the octic construction.


\begin{proof}[Proof of the positive-average bound]
The degree-twenty proposition gives $G/6.16918<J_0(P)\le H(1)-PH(0)$. Apply Proposition~\ref{average} with $a=1/6.16918$ for the right edge. The unitary functional equation and complex conjugation send each zero $\beta+iT$ of the original curve to $1-\beta+iT$. Applying the same right-edge bound to that zero proves the two-sided assertion.
\end{proof}


\section{Fixed-rank zero-free regions over arbitrary number fields}\label{rb:section}
Write $r=[K:\mathbb Q]$ and let $\pi$ be a fixed unitary cuspidal representation of $\mathrm{GL}_n(\mathbb A_K)$, $n\ge2$. The following bounds use the classical zero-free region of Coleman for Hecke characters. They are independent of a fixed zero-free half-plane for the rank-one functions.

\begin{theorem}\label{rb:rank}
For each fixed $\pi$, all its nontrivial standard zeros at sufficiently large $|T|$ satisfy
\[
 \min(\beta,1-\beta)>\frac1{rC_n\log|T|}.
\]
If $\widetilde\pi\simeq\pi\otimes\xi^{-1}$ for a unitary Hecke character $\xi$, one may take $C_2=5.673\cdot2$, $C_3=6.178\cdot3$, and $C_n=6.348n$ in every rank. For a general $\pi$ one may take
\[
 C_n=\min\{6.406n+2.136n^2,\ 6.606n+2.065n^2,\ 6.943n+2.042n^2\}.
\]
The height threshold may depend on the representation and field. No local temperedness is required.
\end{theorem}

\paragraph{The scalar terms.}
For a fixed unitary Hecke character $\eta$ of $K$, put $L=\log|T|$ and fix $u>0$. Uniformly for bounded $v$, one has
\[
 \left|\Re \frac{L'}{L}(1+u/L+i(T+v/L),\eta)\right|=o_\eta(L).
\]
Here and below a fixed norm twist simply translates the ordinate. Coleman, \emph{A zero-free region for the Hecke $L$-functions}, Theorem 2, gives a zero-free width $\delta$ with $\delta L\to\infty$ on the unit neighborhood of the ordinate. For fixed characters the precise Vinogradov--Korobov statement is also given in Petrow, \emph{The Weyl law for algebraic tori}, Lemma 2.12. A possible real exceptional zero causes no difficulty at these ordinates. \href{https://doi.org/10.1112/S0025579300013000}{Coleman}; \href{https://doi.org/10.4171/JEMS/1465}{Petrow}.

Here are details of the deduction. Put $\sigma=1+u/L$, $y=T+v/L-\Im\rho$ and $K_x(y)=x/(x^2+y^2)$. Compare the real Hadamard identities at $\sigma$ and $\sigma+e$, with $e=\sqrt{\delta/L}$. When $|y|\le1$, $x=\sigma-\Re\rho\ge\delta$ and
$K_x(y)/K_{x+e}(y)=1+O(e/\delta)$ uniformly; this follows by placing the ratio between $x/(x+e)$ and $(x+e)/x$. The total of the nonnegative kernels at $\sigma+e$ is $O_\eta(L+1/e)$ by the Hadamard identity and the absolutely convergent Euler derivative. For $|y|>1$, the derivative with respect to $x$ is at most $1/y^2$ in modulus. Summing the fixed-field unit-interval zero counts bounds this part of the difference by $O_\eta(eL)$. The gamma and pole contributions have bounded variation between the two abscissae. Hence the claimed real part is $O_\eta(1/e+(e/\delta)(L+1/e)+eL)=o_\eta(L)$. The same proof holds with the ordinate multiplied by any fixed nonzero integer. Only this real-part estimate is used, including when the scalar coefficient is negative.

\paragraph{The common positive test.}
Let $h$ be the compact nonnegative normalized autocorrelation used in the preceding weighted disk bound, restricted to $[0,1]$, and set
\[
 J(B,P)=B\int_0^1 h(x)\{e^{-Bx}-P\}\,dx .
\]
A prime-power amplifier with normalized target coefficient one, static pole weight $P$, and total shifted disk cost $G$ excludes every limiting scaled right-edge displacement strictly less than $J(B,P)/G$. If $G<CJ(B,P)$ with a strict margin, this excludes displacement at most $1/C$ for all sufficiently large heights. For clarity, the analytic input is the following. Each finite good Euler sum is nonnegative after weighting by the nonnegative autocorrelation at the logarithm of the prime-power norm. Realize this weight by averaging the ordinary disk logarithmic-derivative inequality over the common spectral displacement $v/\log|T|$. At a shifted factor of absolute degree $d$, convexity and the differentiated disk Poisson formula cost $(d/4+o(1))\log|T|$. The left semicircle of the disk contributes the exponent $d/4$ because its cosine-squared integral is $\pi/2$. The projection of each enclosed zero has nonnegative sign. The target contributes its entire reciprocal pole kernel up to a bounded boundary correction. For scalars at a nonzero shifted frequency the preceding classical estimate instead gives $o(\log|T|)$; fixed unshifted poles supply $P$. Integrating the target and static pole kernels uses $\int_{\mathbb R} z/(z^2+v^2)\,d\mu(v)=\int_0^\infty e^{-zx}\widehat\mu(x)\,dx$ for $z>0$. The tails of the spectral average are controlled by absolute convergence at $1+u/\log|T|$, and $u$ tends to zero after taking the high-ordinate limit. This establishes the test. The disk identity itself is classical; see Heath-Brown, \emph{Zero-free regions for Dirichlet $L$-functions}, Lemma 3.2. \href{https://doi.org/10.1112/plms/s3-64.2.265}{Heath-Brown}.

\paragraph{The essentially self-dual amplifier.}
At a fixed unramified prime power let $a$ denote the trace of its Satake parameter, $q$ the underlying norm and $\nu$ the exponent. Choose a unit scalar $z$ with $z^2=\xi(\mathfrak p)^\nu$. The twist relation and unitarity imply $X=a/z\in\mathbb R$. Put $\phi=\arg(z)-\nu T\log q$. For $\alpha>0$ and $0\le s<1$ take the real trace test
\[
 A(\phi)X^2+X\cos\phi+Q(\phi),\qquad
 A(\phi)=\alpha(1+s\cos(2\phi)),\qquad
 Q(\phi)\ge\frac{\cos^2\phi}{4A(\phi)}.
\]
It is nonnegative by completing the square, for every real $X$. The target is the original standard function. The constant-frequency $X^2$ coefficient gives $\alpha$ times the simple pole of $\pi\times\widetilde\pi$. Its frequency-two term is $\alpha s$ times the shift of $(\pi\times\widetilde\pi)\otimes\xi$, of absolute degree $n^2r$. Thus
\[
                  G=\frac{nr}{4}(1+n\alpha s).
\]
The change $z\mapsto-z$ changes both $X$ and $\cos\phi$ in sign. All remaining scalar terms are genuine shifts of powers of $\xi$.

Choose the scalar majorant to be a polynomial in $\cos(2\phi)$, approximating the indicated right side from above with uniform error less than $10^{-6}$. Such a polynomial exists by ordinary uniform polynomial approximation of the continuous function $(1+t)/(8\alpha(1+st))$ on $[-1,1]$. Its nonconstant Fourier coefficients may have either sign by the scalar estimate. For $c>0$ write $s=(1-c^2)/(1+c^2)$. Integrating with normalized circular measure gives the upper pole budget
\[
             P\le P_0+10^{-6},\qquad
             P_0=\alpha+\frac{1+c^2}{8\alpha(1+c)} .
\]
Indeed $\int(1+s\cos(2\phi))^{-1}\,d\phi/(2\pi)=1/\sqrt{1-s^2}$ by the tangent half-angle substitution. When $s=0$ take $Q=\cos^2\phi/(4\alpha)$ exactly, so that $P=P_0$. The required gains follow from the following entirely rational choices.
\[
\begin{array}{c|cccc|c}
 n&\alpha&c&s&B&G/(nrJ(B,P))\text{ is less than}\\\hline
 2&311/1000&19/30&539/1261&89/100&5.673\\
 3&322/1000&7/9&16/65&21/25&6.178\\
 \text{any}&35/99&1&0&8/11&6.348
\end{array}
\]

\paragraph{The unrestricted complex trace.}
For $w=ae^{-i\nu T\log q}$ and $b>0$, take
\[
 \frac{(b+\Re w)^2}{2b}
 =\frac b2+\Re w+\frac{|w|^2}{4b}+\frac{\Re(w^2)}{4b}.
\]
The last term is the shifted $\pi\times\pi$, the third the static $\pi\times\widetilde\pi$. Consequently $P=b/2+1/(4b)$ and $G/r=n/4+n^2/(16b)$. There is no shifted scalar in this test. For $(b,B)=(3/4,18/25),(4/5,71/100),(17/20,69/100)$ the separate ratios $1/(4J)$ and $1/(16bJ)$ are respectively below $(6.406,2.136)$, $(6.606,2.065)$ and $(6.943,2.042)$. These rational bounds and those in the table follow by integrating the ninth-order Taylor minorant of $e^{-Bx}$ against $h$. The Lean module \texttt{Research.RankBump} proves them as inequalities for the actual integral using \texttt{Research.CompactBump}.

The analytic continuation, functional equations, pole orders and convexity used here are the standard Rankin--Selberg theorems of Jacquet, Piatetski-Shapiro and Shalika. A pole of a fixed shifted factor is outside its disk at large $|T|$. At a good place the absolute values of the coefficients of $\pi\times\pi$ are bounded by the nonnegative coefficients of $\pi\times\widetilde\pi$; Cauchy--Schwarz supplies the ordinary standard bound. At the finitely many omitted places the normalized local exponents of a unitary generic standard factor are strictly below $1/2$ and those of a local Rankin--Selberg factor strictly below one. Their logarithmic derivatives therefore stay bounded at the abscissae under consideration. The concluding left-edge assertion follows by the functional-equation symmetry at fixed ordinate. This proves the theorem.


\section{Using an existing zero-free region inside the disk estimate}\label{sec:tilt}
The next step uses a prior strip for a short list of auxiliary functions. It changes the weight at each prime power. A common average still represents this weight, but its measure may have both signs. We therefore work with an identity before applying the disk estimate. The designated weight-one mixed term can be $(1,j_0)$ for any positive symmetric-power index.
\begin{proposition}\label{prop:tilt}
Retain the setting of Proposition~\ref{average}. Assume the ordinary completed Hadamard identity, standard archimedean factors, and $O_F(\log(2+|t|))$ zeros in a unit interval of ordinates for each moving nonscalar factor. Suppose for all its zeros with $|\im\rho-\ell T|\le1$ we have $(1-\re\rho)L\ge d$, for every sufficiently large $|T|$, where $L=\log|T|$ and $d>0$ is fixed. Let $h$ be the even compactly supported autocorrelation of a real nonnegative profile, normalized so that $h(0)=1$. Assume its restriction to $[0,\infty)$ is continuously differentiable and piecewise twice continuously differentiable with integrable second derivative. Write $\mathcal H(z)=\int_0^\infty h(t)e^{-zt}\,dt$ for real $z$. For every $W>0$ for which
\[
       WG<\mathcal H(1-d/W)-P\mathcal H(-d/W)
\]
every zero of the target at sufficiently large $|T|$ satisfies $(1-\beta)L>W$.
\end{proposition}
\begin{proof}
Suppose $w=L(1-\beta)\le W$ along an unbounded sequence. Pass to a subsequence for which $w\to w_*\in[0,W]$. Fix $u>0$ and evaluate at $s_\ell=1+u/L+i\ell T$. On the real line take the even extension of $q(x)=e^{(u+d)x}h(Wx)$ from $x\ge0$. This compact real function is nonnegative. The probability assumption used earlier is unavailable. Instead put $d\nu(v)=\widehat q(v)\,dv/(2\pi)$ using angular frequency. Two integrations by parts on each half-line give $|\widehat q(v)|\ll(1+|v|)^{-2}$, allowing a corner at zero. Fourier inversion gives mass one and cosine transform $q$. Absolute convergence of the prime series and finite total variation now show that the signed average over $s_\ell+iv/L$ represents the Euler sum with multiplier $q(\log(p^\nu)/L)\ge0$. The contribution of bad local factors is $O(1)$.
 
For $\re z>0$ write
\[
 Q_L(z)=L\int_0^\infty q(x)e^{-Lzx}\,dx .
\]
Termwise averaging sends the real zero kernel $\re(1/z)$ to $\re Q_L(z)$. One may interchange the zero sum and this signed average: the unaveraged real kernels are nonnegative, and their sum at $s_\ell+iv/L$ is $O_F(L/u+\log(2+|\ell T+v/L|))$ by the real Hadamard identity and Euler absolute convergence. The bound is integrable against $|\nu|$ for fixed $T$. The same identity at $s_\ell$ may thus be subtracted. Two integrations by parts on the positive half-axis give, uniformly for $\re z>0$,
\[
 \left|Q_L(z)-\frac1z\right|
 \le\frac{|q'(0+)|+\int_0^\infty |q''(x)|\,dx}{L|z|^2}.
\]
For fixed $0<\eta<1/4$, the sum of these differences over $|s_\ell-\rho|\ge\eta$ is $O_{F,\eta}(1)$: for those zeros $|s_\ell-\rho|^{-2}\ll_\eta |2+i\ell T-\rho|^{-2}$, and the real Hadamard identity at $2+i\ell T$ bounds the latter sum by $O_F(L)$. The gamma terms in the signed difference are $o(L)$, since the total mass is one. More explicitly, for each fixed $v$ both gamma derivatives have the same leading multiple of $L$. Their quotients by $L$ have a common integrable majorant $O_F(1+\log(2+|v|))$ by Stirling, uniformly at large $T$ for the fixed shifts. The conductor terms cancel. At the moving elliptic indices there are no poles. If a fixed auxiliary function has a pole at a bounded ordinate, it too lies outside the disk for large $T$ and its contribution to this difference is $o(L)$.

Apply the ordinary, unaveraged disk estimate at $s_\ell$. Keep zeros with $|s_\ell-\rho|<\eta$, dropping the favorable terms from the remaining annulus. Add the signed difference of the full Hadamard identities just obtained. The singular contribution at each retained zero is now $-\re Q_L(s_\ell-\rho)$. The disk correction is at most $\eta/R^2$ per zero, so its sum is $O_F(\eta L)$ with a constant independent of $\eta$. At each of these zeros put $x=L(1-\re\rho)-d\ge0$ and $y=L(\ell T-\im\rho)$. Then
\[
 \re Q_L(s_\ell-\rho)=L\int_0^\infty h(Wt)e^{-xt}\cos(yt)\,dt\ge0 .
\]
For $x>0$ this is half the Fourier transform of the product of the positive-definite functions $h(Wt)$ and $e^{-x|t|}$. Equivalently, it is the average of two positive Poisson kernels under the spectral law of $h(Wt)$. At $x=0$ the result follows by dominated convergence, since $h$ is compactly supported. We can now discard all but the target, which contributes $L\mathcal H((w-d)/W)/W$.

The static poles and signed scalar indices are most simply treated directly in the average. On each bounded $v$-interval Laurent expansion at one and the scalar logarithmic-derivative bound apply uniformly. Outside that interval the absolute Euler bound, divided by $L$, is bounded by $C/u$ times the tail of $|\nu|$. The normalized limiting contribution of the static pole is exactly $P\mathcal H(-d/W)/W$. This follows by integrating $u/(u^2+v^2)$ against $\nu$. Every static regular term and every moving scalar contributes $o(L)$ before normalization. Combining the estimates gives
\[
 \mathcal H((w_*-d)/W)-P\mathcal H(-d/W)\le WG+O(\eta).
\]
The constants in the discarded far terms may grow as $\eta$ decreases, but for each fixed $\eta$ they were divided by $L$ first. Let $\eta\downarrow0$ and use monotonicity of $\mathcal H$ on the entire real axis. Since $w_*\le W$, this contradicts the stated inequality.
\end{proof}


\section{A tilted test at the limiting displacement}\label{rst:section}
Let $K$ have degree $r$ over $\mathbb Q$ and let $\pi$ be a fixed unitary cuspidal representation of $\mathrm{GL}_n(\mathbb A_K)$ with $n\ge2$. Suppose there is a unitary Hecke character $\xi$ with $\widetilde\pi\simeq\pi\otimes\xi^{-1}$.
\begin{theorem}\label{rst:theorem}
At sufficiently large $|T|$ every nontrivial standard zero $\beta+iT$ satisfies
\[
       \min(\beta,1-\beta)>\frac1{5.840\,nr\log|T|}.
\]
\end{theorem}
The gain uses the fact that the same standard function supplies both the zero being tested and every moving nonscalar term. We give the comparison for this case directly.

\begin{proof}
Put $L=\log|T|$, $d=nr$, $\alpha=35/99$, $P=\alpha+1/(8\alpha)$, and $G=d/4$. Use the real trace test of Section~\ref{rb:section} with constant quadratic coefficient. Its Euler coefficient is
\[
 \alpha X^2+X\cos\phi+\frac{\cos^2\phi}{4\alpha}
 =\left(\sqrt{\alpha}X+\frac{\cos\phi}{2\sqrt{\alpha}}\right)^2 .
\]
The only term at a moving nonscalar frequency is $L(s+iT,\pi)$ with coefficient one. The static Rankin--Selberg function and Dedekind zeta function have total pole weight $P$. The nonconstant scalar term is a fixed Hecke function at $2T$, with coefficient $1/(8\alpha)$.

Let $\lambda$ be the lower limit of $(1-\beta)\log|\Im\rho|$ over the zeros with $|\Im\rho|\to\infty$. An infinite lower limit proves the right-edge statement. Otherwise Theorem~\ref{rb:rank} gives $\lambda>0$. Choose $0<a<\lambda$, and a sequence of zeros $\rho_T=\beta_T+iT$ whose scaled displacements tend to $\lambda$. Every zero with $|\Im\rho-T|\le1$ has $(1-\Re\rho)L\ge a$ for large $|T|$. Indeed one may first take the eventual bound with any $a'\in(a,\lambda)$, and $\log|T+O(1)|/L\to1$.

Use the nonnegative even normalized compact autocorrelation $h$ of Section~\ref{rb:section}, with its support $[-1,1]$, and put $H(y)=\int_0^1 h(x)e^{-yx}\,dx$ for real $y$. Fix $u,W>0$. The signed spectral measure $d\nu(v)=\widehat q(v)\,dv/(2\pi)$, with the Fourier transform $\widehat q(v)=\int q(x)e^{-ivx}\,dx$, has characteristic function
\[
                 q(x)=e^{(u+a)|x|} h(Wx).
\]
It is a real even signed measure of mass one. On the positive half-line $q$ and $q'$ are continuous through the support endpoint, with $q''$ integrable there. Fourier inversion and integration by parts on the two half-lines give $|d\nu(v)|\ll_{u,a,W}(1+|v|)^{-2}\,dv$. At $s=1+u/L+iT$ let $D=-L'/L$ for the standard factor. For $\Re z>0$ define
\[
 K_q(z)=L\int_0^\infty q(x)e^{-Lzx}\,dx
       =\int_{\mathbb R}\frac{d\nu(v)}{z+iv/L}.
\]
Two further integrations by parts give
\[
 |K_q(z)-1/z|\le
 \frac{|q'(0+)|+\int_0^\infty|q''(x)|\,dx}{L|z|^2}.
\]
All implied constants below may depend on the fixed $u,a,W$, on $K$ and on the finite list of representations.

Apply the real Hadamard identity at $s+iv/L$, integrate against $\nu$, and subtract its value at $s$. To justify interchange, integrate the nonnegative kernels $\Re(1/(s+iv/L-\rho))$ against $|\nu|$. Their sum is $O_u(L+\log(2+|T+v/L|))$. This follows from the same real Hadamard identity, Stirling, and absolute convergence. The good standard Euler derivative is bounded absolutely by Cauchy--Schwarz from the static Rankin--Selberg and zeta sums at $1+u/L$, each $O_u(L)$. Finitely many omitted local terms are $O(1)$ by the unitary local bounds in Section~\ref{rb:section}. The weighted change in the archimedean terms is $O(1)$. For $|v|\le L|T|/2$ both gamma derivatives have the same logarithmic term with bounded difference; on the complement integrate their logarithmic bound against $(1+|v|)^{-2}\,dv$. Since $\nu$ has mass one, constant terms cancel. For fixed $0<\eta<1/4$, all zeros at distance at least $\eta$ from $s$ have total change $O_\eta(1)$, by the last kernel bound and
\[
  \sum_{|s-\rho|\ge\eta}\frac1{|s-\rho|^2}
  \ll_\eta \sum_{\rho}\frac{\Re(2+iT-\rho)}{|2+iT-\rho|^2}
  \ll_\eta L .
\]
Thus the signed difference in $\Re D$ is the negative sum of $\Re(K_q(s-\rho)-1/(s-\rho))$ over the near zeros, up to $O_\eta(1)$.

Add this equality to the ordinary differentiated disk estimate with convexity cost $(G+o(1))L$. Retain all near zeros in that estimate and discard the others by their nonnegative projected kernels. In a disk of radius between $1/4$ and $1/3$ the retained boundary corrections are $O(\eta L)$, since each is at most $16\eta$ and the unit-disk zero count is $O(L)$. Consequently
\[
 \int \Re D(s+iv/L)\,d\nu(v)
 \le (G+o(1))L+O(\eta L)+O_\eta(1)
      -\sum_{|s-\rho|<\eta}\Re K_q(s-\rho).
\]
For every retained zero, the last kernel is nonnegative. Indeed it equals $L/2$ times the Fourier transform of
$h(Wx)e^{-((1-\Re\rho)L-a)|x|}$ at frequency $L(T-\Im\rho)$. An autocorrelation has nonnegative Fourier transform, and $e^{-b|x|}$ is positive definite for $b\ge0$. For the selected zero the kernel is $LH(((1-\beta_T)L-a)/W)/W$ for all sufficiently large $|T|$.

Finally add the averaged static and scalar terms. The signed average of a static logarithmic derivative has its exact pole kernel $LH(-a/W)/W$ per unit pole weight, up to $O(1)$. For this assertion, remove the pole from the ordinary derivative and treat the difference by the same Hadamard estimate: all remaining zeros have a fixed positive distance from $1$, and the square-reciprocal sum at the static center is bounded. The gamma average is bounded by integrability of $\log(2+|v|)$ against $|\nu|$. The moving Hecke derivative is $o(L)$ uniformly for $|v|\le L|T|/2$ by the fixed-character zero-free estimate proved in Section~\ref{rb:section}. The complement has mass $O((L|T|)^{-1})$, on which the Euler bound is $O_u(L)$. The aggregate incomplete Euler sum is nonnegative because $q\ge0$ multiplies every good prime-power square. The omitted primes contribute $O(1)$ throughout. Divide by $L$, let $|T|\to\infty$ along the selected sequence, and then let $\eta\downarrow0$. We obtain, for every fixed $a\in(0,\lambda)$ and $W>0$,
\[
               H((\lambda-a)/W)-P H(-a/W)\le WG .
\]
For fixed $W$, compact support gives continuity as $a\uparrow\lambda$. Substituting $W=\lambda/\tau$ in the resulting inequality gives
\[
       \lambda d\ \ge\ 4\tau\int_0^1h(x)
                 (1-Pe^{\tau x})\,dx \qquad(\tau>0).
\]

Take $\tau=69/100$. For $0\le y\le1$, the exponential series remainder after degree nine is bounded above by $11/(10\cdot10!)<10^{-6}$ by a geometric series estimate. Integrate the resulting polynomial minorant against $h$. It yields
\[
4\tau\int_0^1h(x)(1-Pe^{\tau x})\,dx
\ \ge\ \frac{159697402124340139245971808241683}
                    {932570038400000000000000000000000}
\ >\ \frac{25}{146}.
\]
The strict margin makes the claimed right-edge statement eventual; the functional equation gives the other edge. The scalar integral inequality is also checked in \texttt{Research.RankSelfTilt}.
\end{proof}


\paragraph{A refinement with a classical cosine profile.}
The same argument with $\alpha=1/\sqrt8$ gives $P=1/\sqrt2$ and $G=nr/4$. Choose $\theta=13221231/10^7$ and put
\[
 P_0=\frac{1-\theta\cot\theta}{\sin^2\theta}.
\]
The compact cosine of Heath-Brown, Lemma 7.5, normalized at the origin and with its scale parameter one, satisfies $H(0)-P_0H(-1)=P_0\cos^2\theta$ \cite{HB}. The interval calculation gives $P_0>P$ and $P_0\cos^2\theta>0.04283241$. Since $H(-1)\ge0$, the limiting argument with $\tau=1$ now gives
\[
 \min(\beta,1-\beta)>\frac{1}{5.83671\,nr\log|T|}
\]
at sufficiently large height for each fixed essentially self-dual cuspidal representation. The earlier rank-two coefficient $5.673\cdot 2r$ is sharper in rank two. The rational certificate above proves $5.840\,nr$ using only the polynomial profile. The script \texttt{diagonal\_prior.py} verifies the trigonometric rounding for this refinement.

\section{Two short certificates for edge improvements}\label{sec:tabtilt}


\paragraph{A tilted amplifier using only the first two moving symmetric powers.}
Let $M_k=\int_0^1h(x)x^k\,dx$, with $h$ the normalized quartic autocorrelation above. Set $d=1079/20000$. The prior bound for the two moving symmetric powers permits this value since $3(6.178)d<1$. For each $(C,b,P,G)$ in the table take $x=b(dC-1)$, $y=bdC$ and
\[
 A=b\sum_{k=0}^{13}M_kx^k/k!,\qquad B=b\biggl(\sum_{k=0}^{13}M_ky^k/k!+
 \frac{M_{14}y^{14}}{14!(1-y/15)}\biggr).
\]
{\small\setlength{\arraycolsep}{3pt}\[
\begin{array}{c|rrrr|cc}
\deg & C & b & P & G & 10^8 A> &10^8B< \\ \hline
8 &2363/400&5083/5000&306229/500000&82457203/186000000&34215621&43613075\\
12&7359/1250&639/625&152703/250000&41475929/93000000&34371819&43870033
\end{array}
\]}
The degree-13 Taylor polynomial minorizes $e^t$ for every real $t$, by Taylor's formula with remainder. Here $0<y<1$ and after term 14 successive terms in $e^{yu}$ have ratio at most $y/15$ on $0\le u\le1$. Since $h\ge0$, the tilted zero term admits the lower bound $A$ and the pole term the upper bound $PB$. Direct use of the outward-rounded rational bounds in the last two columns gives $C(A-PB)-G>67/10^7$ in degree eight and $>57/10^7$ in degree twelve. Thus the tilted comparison gives the respective reciprocal width coefficients $2363/400$ and $7359/1250$. No iteration of the conclusion is needed.

Each coefficient below is an integer in units of $10^{-6}$ multiplying $T_\ell U_j$. We record $(j,c_{\ell j})$ in row $\ell$. The first-two-power restriction refers to terms with $\ell j>0$.
\[
\begin{array}{c|r|l}
\deg&\ell&(j,c_{\ell j})\\\hline
8&0&(0,612458),(2,333139),(4,-65617),(6,22084),(8,-4888)\\
&1&(1,1000000)\\
&2&(0,314944),(2,140213)\\
&4,5,6&(0,-5702)\ ;\ (1,15297)\ ;\ (0,6351)\\
&8&(0,200),(2,138)\\ \hline
12&0&(0,610812),(2,330682),(4,-63903),(6,21213),\\
&& (8,-3618),(10,-1082),(12,1209)\\
&1&(1,1000000)\\
&2&(0,319529),(2,143554)\\
&4,5,6&(0,-10530)\ ;\ (1,13978)\ ;\ (0,6639)\\
&8&(0,289),(2,944)\\
&10,12&(0,-130)\ ;\ (0,-10)
\end{array}
\]
The unchanged homogeneous two-chart recursion terminates with 490 nonnegative coefficient rectangles at depth at most 16 for the octic and 376 at depth at most 15 for the dodecic. Both verifications, together with the moment arithmetic, are reproducible by \texttt{certify\_tilt.py} and the files \texttt{final\_tilt8.json}, \texttt{final\_tilt12.json}. The program constructs the root net from the displayed Chebyshev expansions and evaluates the moments by integrating the degree-nine polynomial $h$ term by term.


\section{A symmetric-square prior and the elliptic strip}\label{sec:final}
The degree-twelve table just displayed needs a prior for only $L_1$ and $L_2$. A direct amplifier for $L_2$ raises that prior. All its moving nonscalar mass is concentrated in the function whose zeros it measures.
\begin{lemma}\label{squaredirect}
For every fixed non-CM elliptic curve $E/\Q$, the lower limit of $(1-\re\rho)\log|\im\rho|$ over the zeros of $L_2$ at increasing absolute ordinate is at least
\[
 \frac{46952503706646856803334747}{553773220000000000000000000}>0.08478 .
\]
\end{lemma}
\begin{proof}
Use the following table in units of $10^{-6}$. It has exactly one mixed term, $T_1U_2$.
\[
\begin{array}{c|l}
\ell &(j,c_{\ell j})\\\hline
0&(0,651818),(2,453334),(4,265488),(6,-81161),(8,13385)\\
1&(0,210217),(2,1000000)\\
2,3,4&(0,265069),\quad(0,-42747),\quad(0,13916)\\
5,6,7,8&(0,-6009),\quad(0,2978),\quad(0,-1527),\quad(0,660)
\end{array}
\]
The homogeneous recurrence of the octic construction certifies its nonnegativity on the square. At depths $3$ through $11$ the terminal counts are $2,2,6,6,16,22,42,30,44$. Their sum is $170$. Thus $P=325909/500000$ and $G=\gamma_2=139/186$. The static higher powers may have signed weights because they are regular at one. As in Proposition~\ref{prop:tilt}, the moving signed scalar terms are treated by the classical zeta logarithmic-derivative bound.

Write $\lambda$ for the lower limit in the statement. If it is infinite we are done. It is positive by the general rank-three bound in Theorem~\ref{rb:rank}. For every $0<d<\lambda$ the unit neighborhoods of the ordinates of a minimizing sequence meet the required prior, since $\log|T+O(1)|/\log|T|\to1$. In the tilted comparison keep the mixed target at index $(1,2)$. Before the contradiction step its conclusion is, for every $W>0$,
\[
 H((\lambda-d)/W)-PH(-d/W)\le WG.
\]
Here $H(t)=\int_0^1h(x)e^{-tx}\,dx$ for the undilated compact polynomial of the preceding sections. Continuity at every real $t$ follows by compact support. Let $d\uparrow\lambda$ with $W$ fixed, then substitute $W=\lambda/\tau$ with $\tau=83/100$. This gives $\lambda G\ge\tau(H(0)-P H(-\tau))$.

Define $U(z)=1+z+z^2/2+z^3/6+z^4/24+z^5/100$. Since $U'-U=z^4(1/120-z/100)\ge0$ for $0\le z\le5/6$, $e^z\le U(z)$ there. Let $M_i=\int_0^1h(x)x^i\,dx$, as computed from the explicit polynomial. Then
\[
\begin{split}
 \overline M(z)&=M_0+M_1z+M_2z^2/2+M_3z^3/6+M_4z^4/24+M_5z^5/100,\\
 (M_0,\ldots,M_5)&=(315/796,\ 423/4378,\ 57/1592,\ 3715/227656,\ 59/6965,\ 1645/341484).
\end{split}
\]
As $h\ge0$, $H(-\tau)\le\overline M(\tau)$. The lower bound in the statement is the exact rational number $\tau(M_0-P\overline M(\tau))/G$.
\end{proof}
\begin{proof}[Proof of Theorem~\ref{main}]
Take $d=4239/50000=0.08478$. The moving nonscalars of the degree-twelve table are copies of $L_1$ at frequencies $1,5$ and of $L_2$ at frequencies $2,8$. The first has limiting displacement at least $1/(2\cdot5.673)>d$ by Theorem~\ref{rb:rank}; the second satisfies Lemma~\ref{squaredirect}. Each strict margin persists on the unit neighborhoods of $\ell T$ after division by $\log|\ell T|$ in place of $\log|T|$, since the ratio tends to one. Only finitely many frequencies occur.
 
Use Proposition~\ref{prop:tilt} with
\[
 W=10000/57693,\quad b=101/100,\quad
 P=152703/250000,\quad G=41475929/93000000 .
\]
Write $r=d/W$ and $M_-(z)=\sum_{i=0}^9 M_i(-z)^i/i!$. For the dilate $h_b$ the right side of the strict-gain test is at least
\[
 b\{M_-(b(1-r))-P\,\overline M(br)\}.
\]
Indeed $b(1-r)>0$ and $br=0.49401246654<5/6$, so the lower and upper Taylor estimates apply on the entire interval. Substitution gives
\[
 b\{M_-(b(1-r))-P\,\overline M(br)\}-WG
     >\frac{77}{100000000}>0 .
\]
This verifies the strict inequality, hence excludes the right edge. The functional equation reflects the same estimate to the left edge.
\end{proof}
For $E$ in arithmetic normalization this places $\re s$ between $1/2+1/(5.7693\log|T|)$ and $3/2-1/(5.7693\log|T|)$. The central line remains $\re s=1$. The proof supplies no numerical value of the starting height for a general $E$.
\begin{corollary}
For every fixed non-CM $E/\Q$ and every $0<c<1/\sqrt{5.7693}$,
\[
 \sum_{\substack{p\le x\\p\nmid N_E}}\frac{a_p(E)}{\sqrt p}\log p
   =O_{E,c}\big(xe^{-c\sqrt{\log x}}\big).
\]
\end{corollary}
\begin{proof}
Use the standard truncated explicit formula for $-L_1'/L_1$ at a half-integer $x$ and a regular ordinate within distance one of $Y=\exp(\sqrt{\log x/5.7693})$ \cite[Chapter 5]{IK}. The Euler coefficients have absolute value at most $2\Lambda(n)$, and the unit-interval zero count is $O_E(\log(2+|t|))$. These bounds give a truncation error $O_E((x/Y+1)\log^C(2xY))$ for an absolute exponent $C$. The zeros with $T_E\le|\im\rho|\le Y+1$ contribute $O_E(x\exp(-\log x/(5.7693\log(Y+1)))\log^2(2Y))$ by the theorem and partial summation of the zero count. The finitely many remaining nontrivial zeros are strictly to the left of one and save a fixed power of $x$. Fixed archimedean terms and the omitted prime powers are absorbed; the latter contribute $O(\sqrt{x}\log^2(2x))$. The half-integer convention gives every sharp prime cutoff. The strict inequality on $c$ absorbs the logarithmic factors.
\end{proof}


\section{An alternative using finitely many real shifts}\label{signed}
We record an independent kernel that yields the finite-height theorem below. It works without a central-line estimate. Fix $u>0$, an integer $m\ge0$, $N=2m+1$, and $w=N/(2N+1)$. For a sequence $F_j$ put
\[
  \mathcal LF=\sum_{i=0}^{m-1}(F_{2i+1}-F_{2i+2})+w F_N,
  \qquad \mathcal DF=F_0-\mathcal LF .
\]
For a zero with coordinates $\beta,y$ relative to the ordinate of evaluation set $f_b(y)=b/(b^2+y^2)$, $K_j=f_{1+u+j-\beta}(y)+f_{u+j+\beta}(y)$. These are the contributions from a reflected pair of zeros.
\begin{theorem}
For all $0\le\beta\le1$ and all $y\in\R$, $\mathcal DK\ge0$. For prime powers the multiplier $\mathcal D(x^j)$ is nonnegative for $0\le x\le1$. On constants the multiplier is $1-w=1/2+1/(8m+6)$.
\end{theorem}
\begin{proof}
At either edge of the strip telescoping gives
\[
    \mathcal DK=f_u(y)+(1-w)f_{u+N}(y)-wf_{u+N+1}(y).
\]
For $a=u+N$, $f_{a+1}(y)/f_a(y)\le(a+1)/a\le(N+1)/N=(1-w)/w$. This proves positivity on the two edges. In the closed unit strip take $z=\beta+iy$ and replace $K_j$ by $(1+u+j-z)^{-1}+(u+j+z)^{-1}$. Its real part is $K_j$. The exponential of minus the corrected sum is holomorphic, bounded in the strip and has modulus at most one on its edges. The three-lines theorem proves the same bound throughout. For $0\le x\le1$, pairing successive powers shows $0\le\mathcal L(x^j)\le x$. Indeed $\mathcal L(x^j)=x(1-x^{2m})/(1+x)+wx^N\le x$. Finally $\mathcal L(1)=w$ and the identity for $1-w$ follows on substitution.
\end{proof}
When the Euler and Hadamard expansions converge, a finite linear combination commutes with either sum. Applying the theorem to each pair allows a selected target to be retained. For a zero on the symmetry axis use half of the paired contribution. Every zero sufficiently near an edge belongs to a distinct pair.

\paragraph{The conductor-eleven local weights.}
For the quartic polynomial form the squared modulus of the trigonometric polynomial with coefficient list $(7,20U_1,23+15U_2,20U_1,7)$. Divide by the coefficient of $T_1 U_1$ in its expansion. At eleven apply the linear functional $U_j\mapsto x^j$ with $x=11^{-\nu/2}$. The matrix with entries $U_i U_j$ maps by the Clebsch--Gordan identity to $\sum_{k=0}^{\min(i,j)} x^{i-k}x^{j-k}$, a Gram matrix. Hence the squared polynomial remains nonnegative. For the sextic case the coefficient list at normalization $10^6$ is, by rows,
\[
\begin{array}{c|l}
\ell&(j,c)\\ \hline
0&(0,462244),(2,444674),(4,47038),(6,1191)\\
1&(1,1000000),(3,293292)\\
2&(0,433762),(2,343537),(6,348)\\
3&(1,179116),(3,1183)\\
4&(0,33882),(6,31)\\
6&(0,439),(2,110).
\end{array}
\]
The accompanying integer certificates verify positivity at good primes and, after the specialization $U_j\mapsto x^j$, at the bad prime. They already work after decreasing the first entry from $462244$ to $462242$.


\section{An explicit strip from height 262144 at conductor eleven}
\label{section:level11-effective}
Write $t=\log|T|$ in this section. All real parts are in unitary normalization. The following bounds concern the undeformed standard function of $11.\mathrm{a1}$.

\begin{theorem}
Every nontrivial zero $\beta+iT$ with $|T|\ge262144$ satisfies
\[
       \min(\beta,1-\beta)>\frac1{12.246\,\log |T|}.
\]
On the first octave $262144\le |T|\le524288$ the unconditional coefficient improves to $11.682$. Using the ordinary Riemann hypothesis for $\zeta$ improves $12.246$ to $12.091$ throughout the same range. The latter coefficient improves to $11.447$ for $t\ge200$, to $11.232$ for $t\ge1000$, and to $11.169$ for $t\ge10000$.
\end{theorem}

\begin{proof}
The standard curve is semistable and non-CM. Newton and Thorne, \emph{Symmetric power functoriality for holomorphic modular forms}, Publ. Math. IHÉS \textbf{134} (2021), 1--116, supply the cuspidal symmetric-power lifts. For their local and archimedean factors we use Martin and Watkins, \emph{Symmetric powers of elliptic curve L-functions}, ANTS VII, Lecture Notes in Comput. Sci. \textbf{4076} (2006), 377--392. We apply the real logarithmic-derivative identity for the completed functions, counting nontrivial zeros with multiplicity and pairing $\beta+i\gamma$ with $1-\beta+i\gamma$. The completed zeta also has poles at $0,1$.

Use the preceding finite alternating comparison, with $0<u\le .02$, $\sigma=1+u$, $N=2m+1$, $w=N/(2N+1)$ and $c=1-w$. Write $a_k=(-1)^k$ for $0\le k<N$ and $a_N=-w$. A useful exact retention holds at the selected ordinate. Put $F_j(b)=1/(u+j+b)+1/(u+j+1-b)$. For $0\le b\le1$ and $j\ge1$, putting $A=u+j$ and $p=b(1-b)$ gives
\[
 F_j(1)-F_j(b)=
  \left(\frac1A+\frac1{A+1}\right)\frac{p}{A(A+1)+p}.
\]
This difference is nonnegative and decreases with $j$. The subtracted functional
\[
 \mathcal L(F)=\sum_{i=0}^{m-1}(F_{2i+1}-F_{2i+2})+w F_N
\]
is nonnegative on such sequences. Consequently $\mathcal L(F(b))\le\mathcal L(F(1))$. Since $1/(u+b)\ge1/(u+1)$, the corrected pair at $b$ retains at least $1/(u+1-b)$. Indeed the remaining endpoint surplus is
\[
      \frac{1-w}{u+N}-\frac{w}{u+N+1}\ge0 .
\]
This argument is purely rational. The off-ordinate positivity for discarding all other pairs is the finite-kernel theorem.

We use two positive elliptic polynomials from the previous manuscript. The quartic one is the normalized square associated to $(7,20U_1,23+15U_2,20U_1,7)$. Its static pole and full shifted half-degree are $P_4=413/900$ and $g_4=9161/3600$. The sixth-degree polynomial is the degree-six Bernstein certificate in the earlier paper, section \emph{Explicit original zero strips over number fields}, in the proof of the $12.326$ bound. The companion coefficient table and its nonnegativity proof are recorded in \texttt{Research.EllipticWeights} and \texttt{Research.BernsteinPositivity}. It has $P_6=115561/250000$, nontrivial shifted half-degree $g_6=2284863/10^6$ and total shifted scalar weight $Z=468083/10^6$. Both polynomials have target weight one. Their Euler sums are nonnegative at all primes, including $11$, with signed-comparison Euler multiplier $1-\mathcal L((p^{-k\nu})_k)>0$. At the single prime of bad reduction this last fact also follows directly. At its $\nu$th power the symmetric-power local traces are $b^j$ with $|b|=11^{-\nu/2}$. They are the $U_j$ moments of the positive probability measure $(1-2bx+b^2)^{-1}(2/\pi)\sqrt{1-x^2}\,dx$ on $[-1,1]$, by the generating function and orthogonality of the Chebyshev polynomials $U_j$. Integration of the corresponding nonnegative good-prime polynomial gives its nonnegative bad-prime value. Set
\begin{align*}
 B_4(t)&=\frac{322}{643}g_4t-.77425,\\
 B_6(t)&=\frac{514}{1027}g_6t-.638+Z\{2\log(t+\log 6)+10/3\}.
\end{align*}
We show that each applicable bound has the form
\[
     \frac1{1+u-\beta}\le \frac{P_i}u+B_i(t)
           \qquad(\beta>1/2).
\]

For clarity we give the archimedean calibration. Let $w_{\ell j}$ be the normalized polynomial coefficients. In the quartic case use $m=160$ and complete every shifted function. In the sixth-degree case take $m=256$ and complete the shifted functions with $j>0$. Estimate the finite logarithmic derivatives at $j=0$, $\ell>0$ below. In the following sums the symbol $\mathcal I$ denotes exactly the completed shifted indices. Put $g=\sum_{\mathcal I}(j+1)w_{\ell j}/2$, $h=\sum_{\ell,j}j w_{\ell j}/2$, and let $u_k=.02$ when $a_k>0$, $0$ otherwise. The static real gamma shifts of index $j$ consist of $(j/2)\bmod2$ and the pairs $b,b+1$, $1\le b\le j/2$, counted with multiplicity. Call this multiset $M_j$. An upper bound for the sum of all regular terms beyond $cgt$ is
\begin{align*}
 C={}&P\left(\frac cN-\frac{w}{N+1+.02}\right)+ch\log11\\
 &+\frac12\sum_j w_{0j}\sum_{\mu\in M_j}
     \left[\sum_{k=0}^N a_k\psi\left(\frac{1+k+u_k+\mu}2\right)
                                     -c\log\pi\right]\\
 &+ c\sum_{(\ell,j)\in\mathcal I}\frac{j+1}{2}w_{\ell j}\log\frac{\ell}{2\pi}
   +E.
\end{align*}
For the static pole, the entire completion $s(s-1)\pi^{-s/2}\Gamma(s/2)\zeta(s)$ contributes $+1/(s-1)+1/s$ to its logarithmic derivative. With $f(x)=1/x$, exact cancellation gives
\[
 \sum_{k=0}^N a_k\{f(u+k)+f(u+k+1)\}
 =f(u)+c f(u+N)-w f(u+N+1).
\]
The upper bound on the last two terms is $c/N-w/(N+1+.02)$. In particular the final term has a negative coefficient. This explains the regular pole contribution in $C$. At eleven, the Weil--Deligne model of the $j$th symmetric power is unramified on inertia with one nilpotent Jordan block of size $j+1$. The Artin conductor exponent is the codimension of the kernel, namely $j$. Every other finite prime is good, so the conductor is $11^j$ and its logarithmic-derivative term is $+(j/2)\log11$. The signed coefficients sum to $c$, giving $ch\log11$. For $j=2d$, the unitary real factors are $\Gamma_{\mathbb R}(s+(d\bmod2))$ and $\Gamma_{\mathbb C}(s+b)$ for $1\le b\le d$. For $j=2d+1$ there are the factors $\Gamma_{\mathbb C}(s+b+1/2)$, $0\le b\le d$. Splitting $\Gamma_{\mathbb C}(z)=\Gamma_{\mathbb R}(z)\Gamma_{\mathbb R}(z+1)$ yields the stated static multiset and gives $0\le\mu\le4$ at every moving index. Each real factor contributes $(\psi((s+\mu)/2)-\log\pi)/2$. Monotonicity of the real digamma gives the indicated one-sided choices of $u_k$. At mobile ordinates the inequality $|\psi(z)-\log z|\le1/|z|$ on the right half-plane gives, with $H=262144$,
\[
 E=g(N+1)\left(\frac2H+\frac{(5.02+N)^2}{2H^2}\right)
           + \mathbf1_{i=4}\left(\sum_{\ell>0}w_{\ell0}\right)\frac{N+1}{H^2}.
\]
Indeed the real shifts have size at most four, $\log(1+x)\le x$, and $\sum|a_k|<N+1$. For the digamma inequality, DLMF 5.9.13 gives $\log z-\psi(z)=\int_0^\infty e^{-zt} f(t)\,dt$ with $f(t)=(1-e^{-t})^{-1}-t^{-1}$. Here $f(0+)=1/2$, $f(+\infty)=1$, and $f'(t)=t^{-2}-(2\sinh(t/2))^{-2}\ge0$. Integration by parts therefore gives $|z(\log z-\psi(z))|\le1$ for $\Re z>0$. When the scalars are completed their shifted pole pair is bounded by the boundary telescope at ordinate $\ell T$, namely $f_u+c f_{u+N}-w f_{u+N+1}$ with $f_a=a/(a^2+\ell^2 T^2)$. This is at most $(N+1)/H^2$ for $N=321$. The accompanying \texttt{level11\_constants.py} computes $C$ by outward-rounded Arb arithmetic at 192 bits. It gives $C<-.77425$ in the quartic calculation and $C<-.638$ in the sixth-degree calculation.

Here is the rank-one estimate supplying the remaining sixth-degree term. Write $L_0=\log(6|T|)>13$, put $X=L_0^2/4$ and take $s=\sigma+i\ell T$ for $\ell=2,4,6$. Shift the absolutely convergent Perron integral for $-\zeta'/\zeta$, with kernel
\[
 J(z)=\frac{(e^4X)^z-X^z}{4z^2}.
\]
Its residue at zero is one. Shift from $\Re z=2$ to the left across $z=0$, $1-s$, $\rho-s$ and $-2k-s$ for $k\ge1$. For the horizontal sides choose ordinates separated from the zeta zeros by $\gg1/\log V$. The local partial-fraction expansion then gives $O(\log^2 V)$ for the logarithmic derivative in each fixed strip. The factor $z^{-2}$ makes those integrals vanish. On left lines avoiding the trivial zeros the functional equation gives logarithmic growth while $X^{\Re z}$ tends to zero. These standard estimates are given in Titchmarsh, \emph{The Theory of the Riemann Zeta-function}, second edition, Chapter IX. Thus the shifted integral is the sum of the stated residues. The prime sum has weights averaging the cutoffs $Xe^v$ over $0\le v\le4$. For $y\ge1$ we use $\sum_{n\le y}\Lambda(n)/n\le\log y+.41$. For an integer $k\ge1000$, the interval $(\lceil k/2\rceil,k]$ contributes at most $k\log 2$ to $\psi(k)$, by the middle binomial coefficient. Iteration gives $\psi(k)\le2k\log2+\log(2k)$. Now expand $\log(k!)$ and bound it by $k(\log k-1)+1+\log k$ to obtain the claimed $.41$ allowance. The error $(2\log k+1+\log2)/k$ is decreasing. The fewer than one thousand remaining integers are verified in the same interval calculation. Thus the polynomial costs at most $\log X+2+.41$.

Assume all nontrivial zeta zeros have real part one half. The sum of $|s-\rho|^{-2}$ is bounded by six times the sum of the real zero kernels at $2+i\ell T$, by direct comparison of their denominators. The latter is at most
\[
 \tfrac12\log(|\ell T|/(2\pi))+.001+.58 \le L_0/2.
\]
Here the real Hadamard identity and the digamma bound give the first two terms. For the last term, direct summation to $1000$ followed by $\int_{1000}^\infty (\log x)x^{-2}dx$ gives $-\zeta'(2)/\zeta(2)<.58$. Thus the total nontrivial residue cost is at most $3(1+e^{-2})/2$. The trivial zeros cost at most $(1+e^{-4})/(2X)$ and the pole costs at most $1/(2H^2)$. Finally $0\le\mathcal L((v^k)_k)\le v$ for $0\le v\le1$ by pairing the alternating powers. Hence the combined finite correction from all $k\ge1$ costs at most $.58$ per scalar weight. The resulting upper cost per scalar weight is $2\log L_0+10/3$, since the absolute constant, after substituting $L_0\ge13$ in the error, is less than $3.319$. This proves both inequalities for $B_i$.

Optimizing at $u=(\sqrt{P_i}-P_i)/B_i(t)$ gives the width $(1-\sqrt{P_i})^2/B_i(t)$. The two $B_i$ increase in $t$ and already at $\log H$ their optimizing values of $u$ lie strictly below $.02$. The ratio $B_4(t)/t$ increases with $t$, and $B_6(t)/t$ decreases. Indeed, for its nonconstant numerator $A=-.638+Z(2\log(t+\log6)+10/3)$ we have $A-t A'\ge-.638+Z(2\log(t+\log6)+10/3-2)>0$. At the split $t=48$, the two reciprocal width coefficients are less than $12.090948$ and $12.090692$, respectively. Apply the quartic bound to the left and the sixth-degree one to the right. Removing the negative intercept from the quartic calculation alone gives $12.246$. At $t=\log524288$ the quartic ratio is below $11.681004$, proving the first-octave statement. Evaluation at $t=200,1000,10000$ gives the other three constants. Reflection gives the opposite edge. The formal theorem \texttt{EffectiveEleven.quartic\_exclusion} checks the transfer from the displayed quartic analytic inequality to the actual numerical exclusion at $|T|\ge262144$. The same module checks the first-octave, two-polynomial and three high-ordinate conclusions. It takes the analytic inequalities as hypotheses. The Archimedean evaluations are separate ball-arithmetic calculations.
\end{proof}


\section{Elliptic curves over other number fields}\label{sec:fields}
\begin{corollary}
Let $K$ be totally real, $r=[K:\Q]$, and let $E/K$ be a fixed modular non-CM elliptic curve. Each nontrivial unitary zero at sufficiently large $|T|$ satisfies
\[
       \min(\beta,1-\beta)>\frac1{7.29\,r\log|T|}.
\]
\end{corollary}
\begin{proof}
Newton--Thorne's symmetric-power theorem for Hilbert modular forms of regular weight gives cuspidal symmetric powers at every index in the sextic table of Section~\ref{signed}. Take that table with first entry $462242$ and normalize $L=r\log|T|$. At good primes the normalized elliptic Frobenius eigenvalues have absolute value one. The scalar logarithmic derivatives have size $o(L)$ by the fixed Hecke estimate in Section~\ref{rb:section}. Each moving positive-degree index $j$ costs $(j+1)/4$ in this normalization by convexity. The pole and cost are $P=462242/10^6$ and $G=2284863/2000000$. The moving functions have rank at most seven and are all self-dual. Theorem~\ref{rst:theorem} therefore supplies the common prior $d=61/2500<1/(7\cdot5.840)$. The table has positive mixed entries.
  
Apply Proposition~\ref{prop:tilt} with $W=100/729$ and $b=41/25$. Its proof uses $L$ only as the normalization of the central growth, the displacements and the prime weights, so it applies verbatim with the present $L$. The rational lower gain calculation gives
\[
 (1/W)b\{M_-(b(1-d/W))-P\overline M(bd/W)\}-G>1/1000,
 \qquad 0<bd/W<5/6 .
\]
This yields the right-edge statement and reflection gives the other edge. Here modularity over $K$ is a hypothesis on $E$. See J. Newton and J. A. Thorne, \emph{Symmetric power functoriality for Hilbert modular forms}, Theorem A, \href{https://api.repository.cam.ac.uk/server/api/core/bitstreams/6bd9d94f-51d3-42e2-9cb1-ffc812189ac5/content}{author manuscript}.
\end{proof}
\paragraph{The CM case.}
If $E/K$ has complex multiplication, the main theorem of complex multiplication identifies its standard $L$-function, up to the normalization shift, with a product of Hecke functions of $K$ or a Hecke function over its quadratic CM endomorphism field. In the latter case induction preserves the $L$-function and in the former the character decomposition gives the product. Apply Coleman's fixed-character high-ordinate zero-free region to the finite list and reflect \cite{Coleman}. In particular, for \emph{every} fixed $A>0$ the unitary edge displacement of every sufficiently high zero exceeds $A/\log|T|$. This abelian conclusion is already available from classical results. It explains the non-CM hypothesis in the constructions using higher symmetric powers. For the CM factorization see J. H. Silverman, \emph{Advanced Topics in the Arithmetic of Elliptic Curves}, Chapter II.


\section{Several zero ordinates in the same prime sum}\label{sec:occupation}
Let $\pi$ be a fixed unitary cuspidal representation on $\mathrm{GL}_n$
over a number field $K$, with $d=n[K:\mathbb Q]$. Count occurrences of
nontrivial zeros with multiplicity. Take $H(T)\ge1$ with
$\log(2+H(T))/\log T\to0$ and $I_T=[T,T+H(T)]$. The ordinary completed
standard function is in unitary normalization. The proof also works at
negative heights by applying it to $\widetilde\pi$.

\paragraph{Theorem.}
For $k\ge2$ and
\[
   0\le A<\frac{4}{d}(1-k^{-1/4})^2,
\]
all sufficiently large $I_T$ contain at most $k-1$ occurrences satisfying
$\beta\ge1-A/\log T$. In particular, at most one occurrence satisfies
$\beta\ge1-1/(9.876d\log T)$. If $\pi$ is essentially self-dual, the
interval may instead be $[T,2T]$, and for every $k\ge1$ the admissible
range increases to $0\le A<4(1-(2k)^{-1/4})^2/d$.
The latter assertion uses the fixed-character estimate of Section~\ref{rb:section}. These upper bounds allow the same zero ordinate to appear
several times when the multiplicity demands it.

\paragraph{Proof.}
Put $L=\log T$, take $k$ candidate zeros $\beta_i+i\gamma_i$, and write
$x_i=(1-\beta_i)L$, $y_i=\gamma_iL$. At an unramified prime power with
norm $q^\nu$, let $a$ be the trace of the Satake parameter raised to its
$\nu$th power. For nonnegative weights $b_i$ with $\sum b_i^2=1$ use
\[
  \tfrac12\left|1+a\sum_i b_i q^{-i\nu\gamma_i}\right|^2 .
\]
The static pole weights of $\zeta_K$ and $L(s,\pi\times\widetilde\pi)$
are each $1/2$. The target copy at $\gamma_i$ has weight $b_i$.
The unordered pair $i<j$ contributes the real logarithmic derivative of
$L(s+i(\gamma_i-\gamma_j),\pi\times\widetilde\pi)$ with weight $b_ib_j$.

Use the nonnegative signed-abscissa combination at $\sigma=1+v/L$ as in
the preceding proof, initially keeping the abscissa formula fixed. Write
$\eta$ for its total gamma multiplier. Values of $\eta$ arbitrarily close
to $1/2$ from above are available. The normalized upper bound has static
pole $1/v$ and target gamma cost $c_\eta\sum b_i$, $c_\eta=\eta d/2$.
All cross gamma costs are $O_{k,\pi}(\log(2+H(T))/L)=o(1)$ by Stirling.
The pole at one of the cross Rankin--Selberg factor contributes at most
$b_ib_j P_v(y_i-y_j)+o(1)$, where $P_u(t)=u/(u^2+t^2)$.
The pole at zero and the shifts of both poles to the right contribute
$O(1/L)$. Bad local factors contribute $O_{k,\pi}(1/L)$ by the strict
unitary inverse-root bounds at the finitely many omitted places.
The errors may depend on the chosen signed formula and on $k,v,\pi$;
they are uniform for $\gamma_i\in I_T$ and for these bounded weights.

Keep all $k$ selected zeros in each of the target copies. After scaling by
$L$, their kernels in copy $i$ have the lower bounds
$P_{v+x_j}(y_i-y_j)+O(1/L)$. The other point of the functional-equation
pair has nonnegative undifferenced kernel, and all its positively or
negatively weighted evaluations strictly to the right of one are $O(1/L)$
for the fixed abscissa formula. For large $T$ the selected zeros near the
right edge and their reflected partners are disjoint occurrence lists.
The remaining complete pairs have nonnegative signed kernel. Thus all the
zero terms used here are kept with their correct multiplicities.

The key point as two ordinates approach one another is the exact inequality
\[
   P_{v+x}(t)\ \ge\ \frac{v}{v+x}P_v(t)\qquad(x\ge0).
\]
Clearing the positive denominators leaves $x(2v+x)t^2\ge0$.
Consequently the pair of cross target contributions dominates its cross
pole whenever
\[
  b_i\frac{v}{v+x_j}+b_j\frac{v}{v+x_i}\ge b_i b_j.
\]
There is no lower bound on $|\gamma_i-\gamma_j|$.

For equal $b_i=1/\sqrt k$ set $r=k^{1/4}$ and $G=c_\eta\sqrt k$.
Assume $0\le x_i\le A<(r-1)^2/G$ and take $v=(r-1)/G$.
Then $v/(v+x_i)\ge1/r\ge b_i$. All cross poles are paid, yielding
\[
        \frac{r^2}{v+A}\le \frac1v+G+o(1).
\]
At $A=(r-1)^2/G$ the two main terms are equal. Below that value the
strict gap contradicts the inequality for all sufficiently large $T$.
Choose the signed combination sufficiently close to the limiting
$c=d/4$ for the desired fixed $A$, and then let $T$ tend to infinity.
The $k=2$ numerical rounding is strict since
$[4(1-2^{-1/4})^2]^{-1}<9.876$.

For an essentially self-dual $\pi$ write
$\widetilde\pi\simeq\pi\otimes\xi^{-1}$. Choose $z$ on the unit circle
with $z^2=\xi(\mathfrak p)^\nu$. At this prime power $X=a/z$ is real.
Put $Y_i=\Re(zq^{-i\nu\gamma_i})$, and use the square
$(X+\sum b_iY_i)^2/2$ with $\sum b_i^2=2$.
The static pole is again one. Between copies $i,j$ the difference scalar
is $\zeta_K$ with coefficient $b_i b_j/2$; the sum scalar is $L(s,\xi)$
with the same coefficient. At $2\gamma_i$ the additional scalar also has
character $\xi$. These statements follow from the real cosine product
identity at the good prime powers. At $\gamma_i+\gamma_j\ge2T$ the Hecke
logarithmic derivative is $o_\pi(L)$ uniformly for $1\le\sigma\le2$,
by the estimate in Section~\ref{rb:section}. In its Hadamard comparison the bound is uniform down to $\sigma=1$, since the near zeros still have $\sigma-\beta\gg\delta$. At differences in $[-T,T]$ the same argument
applies to $\zeta_K$ away from bounded heights; on each fixed compact
set its pole subtracted logarithmic derivative is bounded on $\sigma\ge1$.
Combining the two ranges gives a uniform $o_\pi(L)$ after explicitly
retaining the difference pole. At the other abscissae, strictly to the
right of one, absolute convergence gives bounded scalar logarithmic
derivatives. This proves the same budget with half the cross-pole cost
for the whole dyadic interval. Now $r=(2k)^{1/4}$, $b_i=2/r^2$ and
$G=c_\eta r^2$. The two cross terms cover a pair's pole since
$(b_i+b_j)/r\ge b_i b_j/2$. The same optimization proves the stated
constant.\hfill$\square$

\paragraph{A weighted consequence.}
For arbitrary $\pi$, fixed $v>0$, $c=d/4$ and the subpower window above,
\[
   \limsup_{T\to\infty}
   \sum_{\gamma_\rho\in I_T}
   \left(\frac{v}{v+(1-\beta_\rho)\log T}-cv\right)_+^2\le1 .
\]
Only $x<1/c-v$ can contribute, so the cardinality of the support is
bounded independently of $T$ by the preceding theorem. For fixed signed
multiplier first use $c_\eta$ in the displayed expression. If the square
root $M$ of the sum were greater than $1+\epsilon$, choose $b_i$ equal
to the positive part divided by $M$. Their squares sum to one and
$b_i\le v/(v+x_i)$, so all the cross poles cancel against selected
target terms. The resulting diagonal inequality, multiplied by $v$,
gives $M\le1+o(1)$. For $c_\eta\downarrow c$ the error in the displayed
sum tends to zero uniformly on this bounded support. This proves the
claim. For context, log-free density estimates already imply qualitative
bounds on the number of zeros near the edge; see R. Lemke Oliver and
J. Thorner, \emph{Effective log-free zero density estimates for automorphic
$L$-functions and the Sato--Tate conjecture}, IMRN (2019).
The count here gives a particular integer threshold and a weighted bound
for ordinates that may coincide.

\paragraph{Formal coverage.}
The file \texttt{Research/Occupation.lean} defines the unordered cross
terms and the exact finite budget. Its exported proofs include the
Poisson comparison, compensation for every pair, the weighted Euclidean
bound, exclusion below the optimized threshold, an eventual statement
for tuples of distinct occurrences, and the strict two-occurrence
rounding. The self-dual normalization has its own finite theorem. The
Euler identities, analytic continuation and limiting explicit-formula
estimates above supply the formal budget hypothesis.


\section{Separate degrees for zeros and poles in Brauer descent}
Write $r=[K:\mathbb Q]$. Suppose a potential automorphy extension and a Brauer identity give
\[
 F(s)=U(s)\prod_{i=1}^k f_i(s)^{n_i},\qquad n_i\in\mathbb Z.
\]
Here the $f_i$ are entire functions, none identically zero, over intermediate fields $M_i$, and $U$ is a holomorphic unit in the region under consideration. Suppose each $f_i$ is zero-free in both open-strip edge regions of width $1/(c d_i\log|T|)$ above a height $H_i$, where $d_i=[M_i:\mathbb Q]$ and $c>0$. Put
\[
 D_+=\max(\{1\}\cup\{d_i:n_i>0\}),\qquad
 D_-=\max(\{1\}\cup\{d_i:n_i<0\}).
\]
\paragraph{Theorem.}
The zeros of $F$ avoid the two edge regions of width $1/(cD_+\log|T|)$ once $|T|$ exceeds the largest $H_i$ with $n_i>0$. Its poles avoid those of width $1/(cD_-\log|T|)$ beyond the largest $H_i$ with $n_i<0$. One must also exceed the height where $U$ is a unit. Different Brauer identities may be used to minimize $D_+$ and $D_-$.

\paragraph{Proof.}
At $z$ let $m_i(z)=\operatorname{ord}_z f_i\ge0$. The order of $F$ equals $\sum_i n_i m_i(z)$. In the first pair of regions all terms with $n_i>0$ vanish, leaving a nonpositive sum. The negative-support argument gives a nonnegative sum. For a finite set $P$ let $N_i(P)=\sum_{z\in P}m_i(z)$. Counting multiplicities, the number of zeros in $P$ is at most $\sum_i\max(n_i,0)N_i(P)$, and the number of poles is at most $\sum_i\max(-n_i,0)N_i(P)$. Both inequalities follow term by term from the order identity. For a rational Brauer identity, clearing its positive denominator $N$ gives the same argument for $F^N$. The exact order identity then reads $N\operatorname{ord}_z F=\sum_i n_i m_i(z)$, so both sign tests remain valid and both multiplicity budgets may be divided by $N$.

For a finite correction factor $1-\alpha q^{-s}$ with $q>1$ and $\alpha\ne0$, taking norms places every root on the line $\Re s=(\log|\alpha|)/(\log q)$. For each of these finitely many real parts $b$ lying in $(0,1)$ take $m_b=\min(b,1-b)>0$. Once $|T|\ge\exp(1+1/(cD m_b))$, the edge width is strictly less than $m_b$. The other real parts lie outside the open strip. Every finite product of such factors to signed integer powers is consequently a holomorphic unit on the high edge regions. These assertions are proved for the complex meromorphic germ and its order in \texttt{SignedBrauerDescent.lean}.\hfill$\square$

\paragraph{A sharp $A_5$ example.}
Suppose a suitable potential automorphy extension $F_0/K$ has Galois group $A_5$. The representation-theoretic step gives
\[
 \mathbf1_{A_5}=\operatorname{Ind}_{D_{10}}^{A_5}\mathbf1
              -\operatorname{Ind}_{A_4}^{A_5}\omega,
\]
where $D_{10}$ is dihedral of order ten and $\omega$ is a nontrivial linear character of $A_4/V_4$. Thus the zeros use degree $6r$ and the poles use degree $5r$. The Galois group and the automorphy descent hypotheses for the two constituents are assumptions of this example.

To see the identity, order the conjugacy classes as $1,2,3,5A,5B$. The two induced rows are $(6,2,0,1,1)$ and $(5,1,-1,0,0)$. The centralizer orders are $60,4,3,5,5$. In $D_{10}$ there are five involutions and two elements from each order-five class. In $A_4$ all three involutions lie in the kernel of $\omega$, and its eight elements of order three have character sum $4(\zeta_3+\bar\zeta_3)=-4$. Apply the induced-character formula class by class.

These two degrees are simultaneously optimal among all rational decompositions of the trivial character as signed inductions of linear characters of proper subgroups. For the positive support take the virtual character
\[
 \Psi=\mathbf1+\chi_3+\chi'_3-\chi_4=(3,-1,0,3,3).
\]
Each proper subgroup lies in a conjugate of $A_4,D_{10}$ or $S_3$. On these three maximal subgroups $\Psi$ restricts respectively to the irreducible three-dimensional representation, the sum of one trivial and two sign characters, and the sum of the sign and two-dimensional standard characters. Frobenius reciprocity shows that every allowed induced character has nonnegative pairing with $\Psi$. Only conjugates of $A_4$ have index at most five, and every linear character of $A_4$ pairs to zero with its three-dimensional irreducible. If all positive terms had index at most five, the total pairing would be nonpositive, contradicting $\langle\mathbf1,\Psi\rangle=1$. Some positive coefficient thus requires index at least six. At least one negative coefficient is necessary because each proper monomial induction has degree at least five and trivial multiplicity at most one. Its degree minus trivial multiplicity is strictly positive, whereas this quantity vanishes on $\mathbf1$. This proves the two assertions of optimality within the stated class of character decompositions.\hfill$\square$

\paragraph{Formal inputs for the group calculation.}
The subgroup and class data are recorded in \href{https://www.math.rwth-aachen.de/~Thomas.Breuer/ctbllib/ctbltoc/data/A5.html}{the GAP Character Table Library} and \href{https://people.maths.bris.ac.uk/~matyd/GroupNames/characters.html}{the GroupNames character tables}. The Lean file \texttt{A5BrauerCertificate.lean} checks the fifteen possible input rows obtained by grouping the linear characters of proper subgroups by their sums on elements of each order. The two order-five classes can be averaged since both linear tests have equal weights on them. For the two terms in the identity the individual order-five values already agree. The identification with actual induced characters uses the group facts and the argument above.

\paragraph{The group of order 168.}
For $G=\operatorname{PSL}_2(7)$ there is a second sharp identity,
\[
  \mathbf1_G=\operatorname{Ind}_{B}^{G}\mathbf1-
             \operatorname{Ind}_{S_4}^{G}\mathrm{sgn},
  \qquad B=C_7\rtimes C_3 .
\]
The two degrees are $8$ and $7$. The action on the projective line is doubly transitive, so its character is $1+\chi_7$, where $\chi_7$ is irreducible. The sign induction has dimension seven and has no trivial constituent. The other nontrivial irreducible dimensions are $3,3,6,8$, which cannot sum with nonnegative multiplicities to seven. Hence the sign induction equals $\chi_7$.

The dual witness $\Psi=\mathbf1+\chi_3+\chi'_3+\chi_8-\chi_6$ has values $(9,-3,0,3,2,2)$ on $1,2,3,4,7A,7B$. On each maximal $S_4$ it restricts to three copies of the sign-twisted standard representation. On $B$ it is the sum of the regular representation of $B/C_7$ inflated to $B$ and the six-dimensional augmentation representation from the affine action of $B$ on seven points. Translations of order seven have no fixed point and elements of order three have one, giving the required values $2$ and $0$. The \href{https://www.math.rwth-aachen.de/~Thomas.Breuer/ctbllib/ctbltoc/data/L3(2).html}{maximal subgroups} have precisely these types. The previous dual-witness proof now gives sharp positive index eight and negative index seven. The namespace \texttt{PSL27BrauerCertificate} checks the table arithmetic for 23 subgroup/character types, yielding 22 distinct compressed rows. As in the $A_5$ example the average of the two split classes suffices for both tests, and the two rows in the identity have identical values on those classes.

An independent audit starts with the even permutations of five letters and the invertible binary $3\times3$ matrices acting on their eight vectors. From each subgroup already found, it adjoins a representative of every remaining right coset and takes the subgroup generated. Induction on the size of a generating set shows that this enumerates every subgroup. For each proper subgroup $H$, its linear characters are obtained from each normal $K\le H$ with cyclic $H/K$ by sending a generator to a primitive root of unity. Evaluation modulo the appropriate cyclotomic polynomial checks each order-class sum is an integer, so all primitive choices yield the same compressed row. The audit also checks the two-term identities on the original unmerged classes. For $A_5$ this covers 58 proper subgroups and 158 subgroup/character pairs. For $\mathrm{GL}_3(2)\cong\mathrm{PSL}_2(7)$ it covers 178 proper subgroups and 557 pairs. The executable \texttt{audit\_brauer\_groups.py} uses only the Python standard library.

\paragraph{Theorem for direct-product modularity extensions.}
Let $G=\prod_{i=1}^k G_i$ with $k\ge1$, each $G_i$ either $A_5$ or $\operatorname{PSL}_2(7)$. Assign $(a_i,b_i)=(6,5)$ or $(8,7)$ respectively. Among all rational decompositions of $\mathbf1_G$ by signed inductions of linear characters of \emph{solvable} subgroups, the optimal positive and negative indices are
\[
 A=\prod_i a_i,\qquad B=\max_j\left(b_j\prod_{i\ne j}a_i\right).
\]
One decomposition attains both. Thus if the potential automorphy extension has this Galois group, the sign-sensitive theorem applies with $D_+=rA$ and $D_-=rB$. The conclusion allows subgroups whose coordinates are linked by relations.

\paragraph{Proof.}
Taking the exterior tensor product of the two-term identities gives the upper bounds. The term obtained by choosing only positive constituents has index $A$. Any term with a negative coefficient has at least one negative coordinate. Its index is at most the index obtained by retaining just that negative choice, and singleton negative choices occur.

For the lower bound let $H\le G$ be solvable, $H_i$ its projection and $P=\prod H_i$. Each $H_i$ is proper since $G_i$ is nonsolvable. Write $\Psi_i$ for the corresponding dual witness. Each restriction $\Psi_i|_{H_i}$ is genuine. Every irreducible constituent $V_i$ of that restriction satisfies
\[
  [G_i:H_i]\dim V_i\ \ge a_i.
\]
Indeed the assertion is immediate at index at least $a_i$. The smaller index can only be $b_i$, where all constituents of the restriction have dimension three.

Suppose $\operatorname{Ind}_H^G\chi$ pairs nontrivially with $\boxtimes_i\Psi_i$. By Frobenius reciprocity its pairing is nonnegative, and some irreducible $W=\boxtimes_i V_i$ of $P$ contains $\chi$ upon restriction to $H$. Again by Frobenius reciprocity $W$ occurs in $\operatorname{Ind}_H^P\chi$, so $\prod_i\dim V_i\le[P:H]$. Multiplying by $\prod_i[G_i:H_i]$ proves $[G:H]\ge A$. The global test has pairing one with $\mathbf1_G$, forcing a positive term of index at least $A$.

For any $j$ replace $\Psi_j$ by $\mathrm{reg}_{G_j}-2\mathbf1$. Its pairing with $\mathbf1_{G_j}$ is $-1$, while its restriction to $H_j$ is $[G_j:H_j]\mathrm{reg}_{H_j}-2\mathbf1$, a genuine character since $[G_j:H_j]\ge2$. Every constituent at this coordinate has dimension at least one, giving local cost at least $b_j$. The same argument now forces a negative term of index at least $b_j\prod_{i\ne j}a_i$. Taking the maximum proves the result.\hfill$\square$

In \texttt{TensorBrauerSupport.lean}, the lower-bound theorem takes the constituent-dimension inequality as an explicit hypothesis and multiplies it by the projection indices. Its signed witness argument is over rational coefficients. The two upper-bound lemmas verify the indices of the tensor expansion. The dimension inequality itself also has an elementary proof formalized in \texttt{RepresentationIndex.lean}. Choose a nonzero vector $v$ spanning the $H$-eigenline in $W$. Its $P$-orbit spans an invariant nonzero subspace of the irreducible $W$, hence all of $W$. Representatives of the $[P:H]$ left cosets supply enough orbit vectors to span, because multiplication on the right by $H$ acts on $v$ by a scalar.

\paragraph{Corollary on the extension degree.}
Let the potential automorphy extension $F_0/K$ above have Galois group
$A_5^a\times\operatorname{PSL}_2(7)^b$, with $a+b\ge1$. Assume its solvable-subgroup descents and linear twists supply the signed $L$-function identities and the constituent edge estimates in the theorem. Set
\[
 A=6^a8^b,\qquad
 B=\begin{cases}(5/6)A,&b=0,\\(7/8)A,&b>0.\end{cases}
\]
At sufficiently large $|T|$ the zero-exclusion width is $1/(crA\log|T|)$ and the pole-exclusion width is $1/(crB\log|T|)$. Relative to the common field-degree bound $[F_0:\mathbb Q]=r60^a168^b$, the zero width increases by $10^a21^b$ and the pole width by $(A/B)10^a21^b$. In particular these two factors are $10,12$ for $A_5$ and $21,24$ for $\operatorname{PSL}_2(7)$. Indeed $A$ is the product of the positive indices and $B/A$ is the largest available ratio $b_i/a_i$. The sign-specific height thresholds in the theorem still apply. The optimality statement concerns rational inductions of linear characters of solvable subgroups of this fixed Galois group. The argument assumes the specified automorphy extension exists for the elliptic curve in question and gives no exclusion in the remaining central portion of the critical strip.


\section{Exact evaluation of the compact cosine test}
\label{sec:compact-cosine-resolvent}

Let $P\in(0,1)$ be the normalized static pole coefficient of a nonnegative
amplifier and $G>0$ its weighted logarithmic conductor or central-growth
cost. Suppose its averaged explicit formula gives the following inequality
along a sequence of zeros for which $(1-\beta)\log|T|\to w<\infty$:
\begin{equation}
 \int_0^\infty h(x)\{e^{-(u+w)x}-P e^{-ux}\}\,dx\le G
 \qquad (u>0).
 \label{eq:cosine-premise}
\end{equation}
Here $h$ may be the normalized autocorrelation of a nonzero nonnegative,
compactly supported, continuous $L^2$ function, or a positive dilation
of that autocorrelation. This is the spectral averaging hypothesis
established by averaging the disk inequality. The test functions used below
are precisely the compact cosines of D.~R.~Heath-Brown,
\emph{Zero-free regions for Dirichlet $L$-functions, and the least prime
in an arithmetic progression}, Proc.~Lond.~Math.~Soc.~(3) \textbf{64}
(1992), 265--338, Sect.~7, Lemma~7.4. Take his $\lambda=1$,
$\theta=\pi-\arctan k$, and remove an immaterial scalar. His
variational discussion motivates the shape; the proof of global
maximality for the specified nonnegative $L^1\cap L^2$ class is given
below. For the complementary angle normalization see
K.~Ford, \emph{Zero-free regions for the Riemann zeta function}, in
\emph{Number Theory for the Millennium}, vol.~II (2002), 25--56, and
M.~J.~Mossinghoff, T.~S.~Trudgian and A.~Yang,
\emph{Explicit zero-free regions for the Riemann zeta-function},
Res.~Number Theory \textbf{10} (2024), article~11, Sect.~4.

\paragraph{Theorem.}
Let $k$ be the unique positive solution of
\[
 P=\frac{k^3}{(1+k^2)(k+\pi-\arctan k)} .
\]
Under \eqref{eq:cosine-premise}, every finite limiting width satisfies
\[
 w\ge \frac{1}{G(1+k^2)}.
\]
The conclusion implies strict exclusion with any larger reciprocal
constant and a sufficiently large threshold depending on that constant.

\paragraph{Proof.}
Put $S=(\pi-\arctan k)/k$, $c=\cos(kS)<0$ and
$g(x)=(\cos(kx)-c)\,\mathbf 1_{[-S,S]}(x)$. The cosine decreases from
$1$ to $c$ on $[0,S]$, so $g\ge0$. Define
\[
\begin{gathered} M=\int_{\mathbb R} g(x)\,dx,\quad
 Q=\int_{\mathbb R}g(x)^2\,dx>0,\\
 h(t)=Q^{-1}\int_{\mathbb R} g(x)g(x+t)\,dx,\quad
 H(a)=\int_0^\infty e^{-at}h(t)\,dt.\end{gathered}
\]
The function $h$ is compactly supported, nonnegative and even, with
$h(0)=1$. Its Fourier transform is $|\widehat g|^2/Q$ under the
convention $\widehat g(v)=\int g(x)e^{-ivx}dx$. Thus it is the real
characteristic function of the probability measure with density
$|\widehat g(v)|^2/(2\pi Q)$. Absolute convergence of the double
integrals gives
\[
  H(0)=\frac{M^2}{2Q},\qquad
  H(1)=\frac1{2Q}\int_{-S}^{S}g(x)J(x)\,dx,
 \qquad J(x)=\int_{-S}^{S}e^{-|x-y|}g(y)\,dy .
\]
Set $\lambda=(1+k^2)^{-1}$. For $|x|\le S$, split $J(x)$ at $y=x$.
The two primitives needed on these intervals are
\[
 A(y)=e^y\{\lambda(\cos ky+k\sin ky)-c\},\qquad
 D(y)=e^{-y}\{\lambda(-\cos ky+k\sin ky)+c\}.
\]
Their derivatives are $e^yg(y)$ and $e^{-y}g(y)$ in the interior.
The equation $\sin(kS)=-k c$ gives $A(-S)=D(S)=0$. Consequently
\[
 J(x)=e^{-x}A(x)-e^xD(x)=2\lambda g(x)+2(\lambda-1)c .
\]
Also $M=-2c(1+S)$. The defining equation for $k$ is precisely
$P=(1-\lambda)/(1+S)$. Multiplying the last expression for $J$
by $g$ and integrating therefore gives the exact Rayleigh identity
\begin{equation}
                 H(1)-P H(0)=\lambda .
\label{eq:cosine-rayleigh}
\end{equation}
For existence and uniqueness of $k$, write the right hand side of the
pole equation as the product $k^2/(1+k^2)$ and
$[1+(\pi-\arctan k)/k]^{-1}$. Both factors strictly increase on
$(0,\infty)$, with product limits zero and one.

Dominated convergence allows $u\downarrow0$ in
\eqref{eq:cosine-premise}. For every $W>w$ choose there $h(Wx)$.
After the substitution $t=Wx$, monotonicity of $e^{-at}$ in $a$ yields
$GW\ge H(w/W)-P H(0)\ge H(1)-P H(0)=\lambda$.
Letting $W\downarrow w$ proves the assertion, including the case
$w=0$.\hfill$\square$

\paragraph{Certified examples.}
For the quartic square, the degree-eight table, and the degree-twenty table, bisection with outward rounding gives the following enclosures for $G(1+k^2)$:
\[
\begin{array}{lcc}
\text{table}&\text{lower bound}&\text{upper bound}\\ \hline
\text{quartic square}&6.19293180&6.19293181\\
\text{octic}&6.17619476&6.17619477\\
\text{degree twenty}&6.16908379&6.16908380
\end{array}
\]
The validation script evaluates the defining equation for $k$ at both endpoints of each bracket before evaluating the cost.

\paragraph{Sharpness for autocorrelations of nonnegative profiles.}
For $K(x)=e^{-|x|}$ define
\[
 \mathcal{E}_P(f)=
 \frac{\langle f,K*f\rangle-P\|f\|_1^2}{2\|f\|_2^2},
 \qquad 0\ne f\ge0,\quad f\in L^1(\mathbb R)\cap L^2(\mathbb R).
\]
Then $\sup_f\mathcal{E}_P(f)=\lambda$. This determines the best
common vertical test within this class, for fixed $P$ and $G$.

To see this, maximize the numerator divided by two on the cone
$\|f\|_2\le1$. Its supremum $\lambda_*$ is positive by the preceding
cosine construction. Young's inequality bounds the convolution energy by
$2\|f\|_2^2$. Hence a maximizing sequence may be taken to satisfy
$\|f\|_1\le(2/P)^{1/2}$. The Riesz rearrangement inequality
(Lieb and Loss, \emph{Analysis}, Ch.~3) permits symmetric decreasing
rearrangements. For such functions,
$f(x)\le\|f\|_1/(2|x|)$ when $x\ne0$. Their $L^2$ tails therefore
tend to zero uniformly. Pass to a weakly convergent subsequence in
$L^2$. The nonnegative, symmetric decreasing cone is closed and convex
in $L^2$ and hence weakly closed. Local integrals against bounded
interval indicators converge weakly, so $\|f\|_1\le\liminf\|f_n\|_1$
for the limit. The operator with kernel
$\mathbf1_{[-A,A]}(x)K(x-y)\mathbf1_{[-A,A]}(y)$ is Hilbert--Schmidt
and therefore compact. Its quadratic form converges under this weak
convergence. Young's inequality and the uniform $L^2$ tails allow
$A\to\infty$, giving convergence of the full convolution energy.
Thus the supremum is attained. Homogeneity and positivity force the
maximizer to have norm one. Write $M=\|f\|_1>0$. At this maximizer,
$\lambda_*\le1-PM^2/2<1$.

Set $J=K*f$ and $V=J-PM-2\lambda_* f$. Differentiating the
homogeneous quotient at $f$ along each bounded, nonnegative, compactly
supported variation gives $\int Vv\le0$. It follows that $V\le0$
almost everywhere. The definition of $\lambda_*$ gives
$\int Vf=0$, hence $V=0$ wherever $f>0$, almost everywhere. Thus
$f=(J-PM)_+/(2\lambda_*)$. The convolution $J$ is continuous,
vanishes at infinity and is symmetric decreasing. The last formula
chooses a continuous representative for $f$ and shows that
$\{f>0\}=(-S,S)$ for some $0<S<\infty$.
Since $K''=K-2\delta_0$ in distributions, $J''=J-2f$ classically.
Within $(-S,S)$,
\[
 f''+k^2 f=PM/(2\lambda_*),\qquad k^2=(1-\lambda_*)/\lambda_* .
\]
The even solution has the form $f(x)=A\cos(kx)+b$, where
$b=PM/[2(1-\lambda_*)]>0$. Monotonicity, positivity and
$f(S)=0$ require $A>0$ and $\pi/2<kS\le\pi$.
On $[S,\infty)$, $J(x)=PM e^{-(x-S)}$. Matching first
derivatives at $S$ gives $2\lambda_* Ak\sin(kS)=PM$.
Since $b=-A\cos(kS)$, this is $\sin(kS)=-k\cos(kS)$.
It follows that $kS=\pi-\arctan k$. Integrating $f$ now gives
$M=-2A\cos(kS)(1+S)$; comparison with
$b=PM/[2(1-\lambda_*)]$ yields $P=(1-\lambda_*)/(1+S)$.
Uniqueness of the positive solution to the pole equation proves
$\lambda_*=\lambda$.\hfill$\square$

\paragraph{Verification.}
The Lean module \texttt{CosineResolvent} evaluates the actual interval
integrals defining $M$ and $J$, proves positivity of the cosine profile
on its support and verifies the integrated identity
$\int g(J-PM)=2\lambda\int g^2$. The module \texttt{CosineBounds} also proves, with a degree-21 arctangent
majorant and rational frequencies, the strict numerical ceilings
$6.192932$, $6.176195$, and $6.169084$. The compactness and rearrangement
argument for optimality is a human proof. The numerical brackets are
reproducible with \texttt{cosine\_constants.py}.


\section{What is machine-checked}
The Lean project uses the ordinary real and complex numbers of Mathlib, Lebesgue integrals, Bochner Fourier transforms, meromorphic orders and actual filter limits. Its function-level statements take $F:\C\to\C$ as the continued standard function. The analytic estimate is an explicit parameter of the theorem. In particular, the project does not define automorphic representations and prove the modularity and subconvexity inputs from first principles.

\begin{center}\small
\begin{tabular}{@{}p{.35\linewidth}p{.58\linewidth}@{}}\toprule
Source module or group & Mathematical statement checked\\\midrule
\texttt{PoissonJensenCore}, \texttt{BlaschkeDisk} & The horizontal derivative of the harmonic Poisson integral, the finite Blaschke product identity and the projected zero kernels.\\
\texttt{SpectralWindow}, \texttt{BumpSpectral}, \texttt{BumpLimit} & The actual spectral probability of the compact polynomial and the integral giving the pole and target contributions.\\
\texttt{CompactBump}, \texttt{TiltedGain}, \texttt{TiltedWindow} & Autocorrelation identities, signed-test zero positivity, and upper and lower exponential moment bounds.\\
\texttt{SquarePriorGain}, \texttt{BootstrapElliptic} & The symmetric-square limiting gain and the elliptic conclusion from separate bounds for the four moving evaluations and the static sum.\\
\texttt{SignedHadamard}, \texttt{SignedPoisson}, \texttt{KernelRemainder} & Interchange with a signed measure, the real zero-kernel identity, and the twice-integrated exponential remainder.\\\\
\texttt{RankBump}, \texttt{RankSelfTilt} & Real integral inequalities for the general-rank constants.\\
\texttt{MultiShift}, \texttt{StripMinimum}, \texttt{SignedKernel}, \texttt{HadamardBudget} & Full-strip positivity and extraction of a target from convergent Euler and zero sums.\\
\texttt{ExactTarget}, \texttt{EffectiveEleven} & On-ordinate retention and the stated conductor-eleven exclusions from the numerical analytic budgets.\\
\texttt{Occupation}, \texttt{OccupationEuler} & Poisson-kernel compensation for several zeros and the local Euler identities.\\
\texttt{SignedBrauerDescent}, \texttt{A5BrauerCertificate}, \texttt{RepresentationIndex} & Meromorphic-order accounting, character tables, and the finite-index eigenline dimension argument.\\
\bottomrule\end{tabular}\end{center}
The matching README identifies certificate files and verification commands. A dependency on an analytic estimate expressed as a theorem hypothesis survives compilation. Checking a numerical theorem about that estimate does not supply an instance for a particular curve. The preceding sections establish those estimates on paper using the cited arithmetic results. None of the formal theorems asserts the extended Riemann hypothesis.
\begin{thebibliography}{20}
\bibitem{BCDT} C. Breuil, B. Conrad, F. Diamond and R. Taylor, \emph{On the modularity of elliptic curves over $\Q$}, J. Amer. Math. Soc. 14 (2001), 843--939. \href{https://doi.org/10.1090/S0894-0347-01-00370-8}{DOI}.
\bibitem{NT} J. Newton and J. A. Thorne, \emph{Symmetric power functoriality for holomorphic modular forms, II}, Publ. Math. IH\'ES 134 (2021), 117--152. \href{https://doi.org/10.1007/s10240-021-00126-4}{DOI}.
\bibitem{BMN} A. R. Booker, M. B. Milinovich and N. Ng, \emph{Subconvexity for modular form $L$-functions in the $t$ aspect}, Adv. Math. 341 (2019), 299--335. \href{https://doi.org/10.1016/j.aim.2018.10.037}{DOI}.
\bibitem{Nelson} P. D. Nelson, \emph{Bounds for standard $L$-functions}, Theorem 1.1. \href{https://arxiv.org/abs/2109.15230}{arXiv:2109.15230}.
\bibitem{HB} D. R. Heath-Brown, \emph{Zero-free regions for Dirichlet $L$-functions, and the least prime in an arithmetic progression}, Proc. London Math. Soc. 64 (1992), 265--338, particularly Lemmas 3.2, 7.4 and 7.5. \href{https://doi.org/10.1112/plms/s3-64.2.265}{DOI}.
\bibitem{JS} H. Jacquet and J. Shalika, \emph{On Euler products and the classification of automorphic representations. I, II}, Amer. J. Math. 103 (1981), 499--558 and 777--815.
\bibitem{RS} H. Jacquet, I. I. Piatetski-Shapiro and J. Shalika, \emph{Rankin--Selberg convolutions}, Amer. J. Math. 105 (1983), 367--464.
\bibitem{Titch} E. C. Titchmarsh, \emph{The Theory of the Riemann Zeta-function}, second edition, revised by D. R. Heath-Brown, Oxford, 1986, Chapters III and VI.
\bibitem{Coleman} M. D. Coleman, \emph{A zero-free region for the Hecke $L$-functions}, Mathematika 37 (1990), 287--304. \href{https://doi.org/10.1112/S0025579300013000}{DOI}.
\bibitem{IK} H. Iwaniec and E. Kowalski, \emph{Analytic Number Theory}, American Mathematical Society Colloquium Publications 53 (2004), Chapter 5.
\end{thebibliography}
\end{document}
